% ============================================================
%  GUÍA TEMÁTICA — MATRICES
%  Curso: 5to Año — Matemática
%  Autor: Ing. Gabriel Astudillo
%
%  CONTENIDO
%    §1 Definición, filas y columnas
%    §2 Orden de las matrices
%    §3 Operaciones con matrices
%        3.1 Multiplicación por un escalar
%        3.2 Suma y resta
%        3.3 Sistemas con matrices (2x2, 3x3 y casos 2x3 / 3x2)
%        3.4 Multiplicación de matrices
%    §4 Matrices equivalentes
%    §5 Operaciones que producen matrices equivalentes
%    §6 Sistemas por Gauss y Gauss-Jordan
%    §7 Determinantes (2x2, 3x3 por Sarrus, 4x4 por cofactores)
%    §8 Regla de Cramer
%    §9 Matriz inversa (por Gauss-Jordan sobre [A|I])
%    §10 Formulario resumen
%    §11 Banco de 88 ejercicios
%    Clave de respuestas comentada (entre \ifmwclaves ... \fi)
% ============================================================

\documentclass[12pt, letterpaper]{article}

% ── Paquetes ─────────────────────────────────────────────────
\usepackage{mathwizards-guia}
\mwguia{Guía Temática — Matrices}{Ing. Gabriel Astudillo $\cdot$ 5to Año — Matemática}
\providecommand{\sen}{\sin}
\setlength{\emergencystretch}{2em}

\usetikzlibrary{matrix, fit, backgrounds, decorations.pathreplacing}

% ============================================================
%  MACROS DE MATRICES
% ============================================================

% Matriz entre corchetes y determinante entre barras verticales.
% Uso: \mwmat{1 & 2 \\ 3 & 4}   \mwdet{1 & 2 \\ 3 & 4}
\newcommand{\mwmat}[1]{\ensuremath{\begin{bmatrix}#1\end{bmatrix}}}
\newcommand{\mwdet}[1]{\ensuremath{\begin{vmatrix}#1\end{vmatrix}}}

% Matriz aumentada: \mwaug{cc|c}{1 & 2 & 3 \\ 4 & 5 & 6}
\newcommand{\mwaug}[2]{\ensuremath{\left[\begin{array}{@{}#1@{}}#2\end{array}\right]}}

% ============================================================
%  ESTILOS TikZ PARA LAS FIGURAS
% ============================================================
\tikzset{
  mwcelda/.style      = {draw=mwprimario!45, fill=white, minimum width=1.18cm,
                         minimum height=0.82cm, inner sep=1pt, font=\small,
                         anchor=center},
  mwceldafila/.style  = {mwcelda, fill=mwnaranja!20, draw=mwnaranja},
  mwceldacol/.style   = {mwcelda, draw=mwmagenta, line width=1pt},
  mwceldagris/.style  = {draw=mwgris!45, minimum width=1.18cm,
                         minimum height=0.82cm, inner sep=1pt, font=\small,
                         anchor=center, text=mwgris},
  mwceldaborrada/.style = {mwceldagris, fill=mwgris!18, text=mwgris!55},
  mwceldapiv/.style   = {mwcelda, fill=mwvioleta!20, draw=mwvioleta,
                         line width=1pt},
  mwpaso/.style       = {draw=mwprimario, fill=mwprimario!8,
                         rounded corners=4pt, minimum width=2.75cm,
                         minimum height=0.8cm, font=\scriptsize, align=center},
  mwflecha/.style     = {-{Stealth[length=2.2mm]}, mwprimario, thick},
}

% ============================================================
%  BANCO DE EJERCICIOS CON NUMERACIÓN CONTINUA
% ============================================================
\newcommand{\mwcheck}[2]{%
  \ifnum\value{enumi}=#2\relax
    \typeout{MW: #1 correcto (#2)}%
  \else
    \GenericWarning{}{MW Warning: #1 termina en \the\value{enumi}; se esperaba #2}%
  \fi}
\newcommand{\mwsecname}{}
\newcommand{\mwsecexp}{0}
\newcommand{\mwsecstart}{0}
% Uso: \mwbloquehead{título}{último nº del bloque}{último nº del bloque anterior}
\newcommand{\mwbloquehead}[3]{%
  \renewcommand{\mwsecname}{#1}%
  \renewcommand{\mwsecexp}{#2}%
  \renewcommand{\mwsecstart}{#3}}
\newenvironment{mwbloque}{%
  \subsection*{\mwsecname}%
  \begin{enumerate}[leftmargin=*, label=\textbf{\arabic*.}]%
  \setcounter{enumi}{\mwsecstart}%
}{%
  \mwcheck{\mwsecname}{\mwsecexp}%
  \end{enumerate}}

\begin{document}
% ============================================================

\renewcommand{\arraystretch}{1.15}

% Un poco de aire después de las estrellas de dificultad
\renewcommand{\facil}{{\color{mwnaranja}$\bigstar\!\!$}\hspace{2pt}}
\renewcommand{\medio}{{\color{mwrojonaranja}$\bigstar\!\bigstar\!\!$}\hspace{2pt}}
\renewcommand{\dificil}{{\color{mwprimario}$\bigstar\!\bigstar\!\bigstar\!\!$}\hspace{2pt}}

\mwportada{Matrices}{Guía temática completa}{De la definición a la inversa, paso a paso}

\vspace{4pt}
\begin{cuadroinfo}
  \small
  \textbf{Presentación.} Una \textbf{matriz} es una de esas ideas que
  aparecen en todas partes: en una planilla de notas, en la lista de precios
  de un inventario, en la programación de una computadora y, sobre todo, en
  el álgebra lineal. Su fuerza está en que permite \emph{ordenar} mucha
  información y \emph{operar} con ella mediante reglas fijas.
  \\[4pt]
  \textbf{Qué vas a aprender:}
  \begin{itemize}[leftmargin=1.2em, nosep, topsep=2pt]
    \item definir una matriz y ubicar cada elemento por fila y columna;
    \item determinar el \textbf{orden} de una matriz y reconocer las matrices
          especiales;
    \item multiplicar por un escalar, sumar, restar y multiplicar matrices,
          sabiendo \emph{cuándo} cada operación es posible;
    \item escribir un sistema de ecuaciones en forma matricial y resolverlo
          por \textbf{Gauss}, \textbf{Gauss-Jordan} y \textbf{Cramer};
    \item calcular determinantes de $2\times2$, de $3\times3$ por
          \textbf{Sarrus} y de $4\times4$ o mayores por \textbf{cofactores};
    \item reconocer y producir \textbf{matrices equivalentes} mediante
          operaciones elementales de fila;
    \item calcular la \textbf{matriz inversa} reduciendo $[A\mid I]$ y usarla
          para resolver sistemas.
  \end{itemize}
  \par\medskip
  \textbf{Cómo usar la guía.} Lee cada sección con lápiz en mano y \emph{reproduce
  cada ejemplo} antes de mirar la solución. Los recuadros rojos son teoría
  (definiciones y propiedades), los magentas son ejemplos resueltos, los
  naranjas son procedimientos y los rojo-anaranjados son advertencias sobre
  errores clásicos. Al final, el banco de \textbf{88 ejercicios} te permite
  comprobar lo aprendido con la clave comentada.
\end{cuadroinfo}

\begin{cuadroteoria}
  \small
  \textbf{El mapa de la guía}\\[6pt]
  \begin{center}
  \begin{tikzpicture}[node distance=0.35cm and 0.35cm]
    \node (n1) [mwpaso] {Definición\\filas y columnas};
    \node (n2) [mwpaso, right=of n1] {Orden de\\la matriz};
    \node (n3) [mwpaso, right=of n2] {Operaciones\\básicas};
    \node (n4) [mwpaso, below=0.85cm of n1] {Sistemas\\$A\mathbf{x}=\mathbf{b}$};
    \node (n5) [mwpaso, right=of n4] {Matrices\\equivalentes};
    \node (n6) [mwpaso, right=of n5] {Gauss y\\Gauss-Jordan};
    \node (n7) [mwpaso, below=0.85cm of n4] {Determinantes\\Sarrus y cofactores};
    \node (n8) [mwpaso, right=of n7] {Regla de\\Cramer};
    \node (n9) [mwpaso, right=of n8] {Matriz\\inversa};
    \foreach \a/\b in {n1/n2, n2/n3, n4/n5, n5/n6, n7/n8, n8/n9}
      \draw[mwflecha] (\a) -- (\b);
    \draw[mwflecha] (n3.south) -- ++(0,-0.45) -| (n4.north);
    \draw[mwflecha] (n6.south) -- ++(0,-0.45) -| (n7.north);
  \end{tikzpicture}
  \end{center}
\end{cuadroteoria}

% ============================================================
\section{Definición, filas y columnas}
% ============================================================

\begin{cuadroteoria}
  \small
  \textbf{Definición.} Una \textbf{matriz} es un arreglo rectangular de
  números ordenados en \textbf{filas} (horizontales) y \textbf{columnas}
  (verticales). Los números del arreglo se llaman \textbf{elementos} o
  \textbf{entradas} de la matriz.
  \[
    A=\mwmat{2 & -1 & 3\\ 0 & 5 & -4}
  \]
  Esta matriz tiene $2$ filas y $3$ columnas: se dice que es de orden
  $2\times3$ y que tiene $2\cdot3=6$ elementos.
\end{cuadroteoria}

Las filas se leen de \textbf{arriba hacia abajo} y las columnas de
\textbf{izquierda a derecha}. En el ejemplo anterior:

\begin{center}
\small
\begin{tabular}{@{}ll@{}}
  \toprule
  \textbf{Concepto} & \textbf{En la matriz $A$} \\
  \midrule
  Primera fila    & $(2,\;-1,\;3)$ \\
  Segunda fila    & $(0,\;5,\;-4)$ \\
  Primera columna & $(2,\;0)$ \\
  Segunda columna & $(-1,\;5)$ \\
  Tercera columna & $(3,\;-4)$ \\
  \bottomrule
\end{tabular}
\end{center}

\begin{center}
\begin{tikzpicture}[baseline=(M.center)]
  \matrix (M) [matrix of math nodes, nodes={mwcelda},
               row sep=-\pgflinewidth, column sep=-\pgflinewidth,
               row 1/.style={nodes={mwceldafila}},
               column 2/.style={nodes={draw=mwmagenta, line width=1pt}}] {
    a_{11} & a_{12} & a_{13}\\
    a_{21} & a_{22} & a_{23}\\
  };
  \node[left=0.45cm of M-1-1, font=\scriptsize, mwnaranja] {fila 1};
  \node[left=0.45cm of M-2-1, font=\scriptsize, mwprimario] {fila 2};
  \node[above=0.3cm of M-1-2, font=\scriptsize, mwmagenta] {columna 2};
  \draw[decorate, decoration={brace, amplitude=4pt, mirror}, mwmagenta]
    (M-2-2.south west) -- (M-2-2.south east);
  \node[below=0.35cm of M-2-2, font=\scriptsize, mwmagenta] {columna};
\end{tikzpicture}
\end{center}

\subsection*{Notación de los elementos}

Cada elemento se identifica con \textbf{dos subíndices}:
\[
  a_{ij} \quad\Longleftrightarrow\quad
  \begin{array}{l}
    i=\text{número de \textbf{fila} (de arriba hacia abajo)},\\[2pt]
    j=\text{número de \textbf{columna} (de izquierda a derecha)}.
  \end{array}
\]
Así, en la matriz $A$ de arriba:
\[
  a_{11}=2,\qquad a_{12}=-1,\qquad a_{13}=3,\qquad
  a_{21}=0,\qquad a_{22}=5,\qquad a_{23}=-4.
\]

\begin{cuadroadvertencia}
  \small
  \textbf{El orden de los subíndices no es negociable.} $a_{12}$ es el
  elemento de la fila $1$ y la columna $2$; $a_{21}$ es el de la fila $2$ y
  la columna $1$. En una matriz no cuadrada son números distintos (aquí
  $a_{12}=-1$ y $a_{21}=0$). Escribe siempre \emph{primero la fila, después
  la columna}: es la misma idea que usar $(x,y)$ para un punto, donde
  primero va la horizontal.
\end{cuadroadvertencia}

\subsection*{Igualdad de matrices}

\begin{cuadroteoria}
  \small
  Dos matrices $A$ y $B$ son \textbf{iguales} ($A=B$) si y solo si:
  \begin{enumerate}[leftmargin=1.4em, nosep]
    \item tienen el \textbf{mismo orden}, y
    \item cada elemento de $A$ es igual al elemento que ocupa la
          \textbf{misma posición} en $B$: $a_{ij}=b_{ij}$ para todo $i,j$.
  \end{enumerate}
\end{cuadroteoria}

\begin{cuadrosolucion}
  \small
  \textbf{Ejemplo.} Determinar $x$ e $y$ para que
  \[
    \mwmat{x+1 & 3\\ 4 & 2y}=\mwmat{5 & 3\\ 4 & 10}.
  \]
  Ambas son de orden $2\times2$. Igualando posición por posición:
  \[
    x+1=5\;\Rightarrow\;x=4,
    \qquad 2y=10\;\Rightarrow\;y=5.
  \]
  \textbf{Comprobación:} con $x=4$ y $y=5$ la primera matriz es
  $\mwmat{5&3\\4&10}$, idéntica a la segunda.
\end{cuadrosolucion}

\subsection*{Matrices especiales}

Estas matrices aparecerán una y otra vez en el resto de la guía; conviene
tenerlas identificadas.

\begin{center}
\small
\begin{tabular}{@{}p{0.26\linewidth} p{0.30\linewidth} p{0.36\linewidth}@{}}
  \toprule
  \textbf{Nombre} & \textbf{Forma} & \textbf{Qué la caracteriza} \\
  \midrule
  Matriz nula $O$ & $\mwmat{0&0\\0&0}$ &
    Todos sus elementos son $0$. \\[6pt]
  Matriz identidad $I_n$ & $\mwmat{1&0\\0&1}$ &
    Cuadrada; $1$ en la diagonal principal y $0$ fuera de ella. \\[6pt]
  Matriz diagonal & $\mwmat{3&0\\0&-2}$ &
    Cuadrada; solo puede tener elementos no nulos en la diagonal. \\[6pt]
  Triangular superior & $\mwmat{2&5\\0&-1}$ &
    Cuadrada; todos los elementos \emph{debajo} de la diagonal son $0$. \\[6pt]
  Triangular inferior & $\mwmat{2&0\\5&-1}$ &
    Cuadrada; todos los elementos \emph{encima} de la diagonal son $0$. \\[6pt]
  Matriz fila & $\mwmat{4 & -1 & 7}$ &
    Tiene una sola fila ($1\times n$). \\[6pt]
  Matriz columna & $\mwmat{4\\-1\\7}$ &
    Tiene una sola columna ($m\times1$). \\
  \bottomrule
\end{tabular}
\end{center}

\begin{cuadroextra}
  \small
  \textbf{La diagonal principal} de una matriz cuadrada son los elementos
  $a_{11},a_{22},\dots,a_{nn}$, es decir, los que van de la esquina superior
  izquierda a la inferior derecha. La \textbf{diagonal secundaria} va de la
  esquina superior derecha a la inferior izquierda. Solo las matrices
  \emph{cuadradas} tienen diagonal principal en sentido estricto.
\end{cuadroextra}

\begin{cuadrosolucion}
  \small
  \textbf{Ejemplo 1.} Para
  $B=\mwmat{1 & 0 & -2\\ 3 & 4 & 5}$, indicar su orden, sus filas, sus
  columnas y los elementos $b_{13}$, $b_{22}$ y $b_{31}$.

  \medskip
  \textbf{Solución.} $B$ tiene $2$ filas y $3$ columnas: orden $2\times3$.
  \[
    \text{Filas: } (1,0,-2)\ \text{y}\ (3,4,5);
    \qquad
    \text{Columnas: } \mwmat{1\\3},\ \mwmat{0\\4},\ \mwmat{-2\\5}.
  \]
  \[
    b_{13}=-2,\qquad b_{22}=4,\qquad b_{31}\ \text{no existe (solo hay 2 filas).}
  \]
\end{cuadrosolucion}

\begin{cuadrosolucion}
  \small
  \textbf{Ejemplo 2.} Construir la matriz $C=(c_{ij})$ de orden $3\times2$
  definida por $c_{ij}=2i-j$.

  \medskip
  \textbf{Solución.} Se calcula cada posición sustituyendo $i$ (fila) y $j$
  (columna):
  \[
    \begin{array}{ll}
      c_{11}=2(1)-1=1, & c_{12}=2(1)-2=0,\\[2pt]
      c_{21}=2(2)-1=3, & c_{22}=2(2)-2=2,\\[2pt]
      c_{31}=2(3)-1=5, & c_{32}=2(3)-2=4.
    \end{array}
    \qquad
    C=\mwmat{1 & 0\\ 3 & 2\\ 5 & 4}.
  \]
\end{cuadrosolucion}

\begin{cuadrosolucion}
  \small
  \textbf{Ejemplo 3.} ¿Son iguales
  $P=\mwmat{2 & -1\\ 0 & 4}$ y $Q=\mwmat{2 & -1 & 0\\ 0 & 4 & 0}$?

  \medskip
  \textbf{Solución.} No. Aunque los números «se parecen», $P$ es de orden
  $2\times2$ y $Q$ es de orden $2\times3$. La igualdad exige, antes que nada,
  el mismo orden.
\end{cuadrosolucion}

% ============================================================
\section{Orden de las matrices}
% ============================================================

\begin{cuadroteoria}
  \small
  \textbf{Definición.} El \textbf{orden} (o \textbf{dimensión}, o
  \textbf{tamaño}) de una matriz es el par
  \[
    m\times n
    \qquad
    \begin{array}{l}
      m=\text{número de filas},\\[2pt]
      n=\text{número de columnas}.
    \end{array}
  \]
  Se lee «$m$ por $n$» y \textbf{siempre se escribe primero el número de
  filas}. Una matriz de orden $m\times n$ tiene exactamente $m\cdot n$
  elementos.
\end{cuadroteoria}

\begin{center}
\begin{tikzpicture}[baseline=(M.center)]
  \matrix (M) [matrix of math nodes, nodes={mwcelda},
               row sep=-\pgflinewidth, column sep=-\pgflinewidth] {
    a_{11} & a_{12} & a_{13}\\
    a_{21} & a_{22} & a_{23}\\
  };
  \draw[decorate, decoration={brace, amplitude=5pt}, mwnaranja]
    (M-1-3.north east) -- (M-2-3.south east);
  \node[right=0.35cm of M-2-3, font=\small, mwnaranja, align=left]
    {$m=2$ filas};
  \draw[decorate, decoration={brace, amplitude=5pt, mirror}, mwmagenta]
    (M-2-1.south west) -- (M-2-3.south east);
  \node[below=0.4cm of M-2-2, font=\small, mwmagenta]
    {$n=3$ columnas};
  \node[below=1.05cm of M-2-2, font=\small, mwgris]
    {orden $2\times3$ $\;\Rightarrow\;$ $2\cdot3=6$ elementos};
\end{tikzpicture}
\end{center}

\subsection*{Clasificación según el orden}

\begin{center}
\small
\begin{tabular}{@{}lll@{}}
  \toprule
  \textbf{Tipo} & \textbf{Condición} & \textbf{Ejemplos de orden} \\
  \midrule
  Matriz cuadrada     & $m=n$      & $2\times2$, $3\times3$, $n\times n$ \\
  Matriz rectangular  & $m\neq n$  & $2\times3$, $3\times2$, $1\times4$ \\
  Matriz fila         & $m=1$      & $1\times2$, $1\times5$ \\
  Matriz columna      & $n=1$      & $2\times1$, $4\times1$ \\
  \bottomrule
\end{tabular}
\end{center}

\begin{cuadroteoria}
  \small
  \textbf{Consecuencias inmediatas del orden:}
  \begin{enumerate}[leftmargin=1.4em, nosep]
    \item Solo se pueden \textbf{sumar o restar} matrices del \emph{mismo}
          orden.
    \item Solo se pueden \textbf{multiplicar} $A$ y $B$ si el número de
          columnas de $A$ coincide con el número de filas de $B$.
    \item Una matriz cuadrada $n\times n$ tiene $n^2$ elementos y es la
          única a la que se le puede calcular determinante en el sentido
          clásico.
    \item El orden \textbf{no} dice cuántos elementos distintos hay: dice
          cuántas posiciones existen.
  \end{enumerate}
\end{cuadroteoria}

\begin{cuadrosolucion}
  \small
  \textbf{Ejemplo 1.} Indicar el orden y el número de elementos de cada
  matriz.
  \[
    A=\mwmat{1&2&3\\4&5&6},\qquad
    B=\mwmat{1&2\\3&4\\5&6},\qquad
    C=\mwmat{1&2&3\\4&5&6\\7&8&9},\qquad
    D=\mwmat{7&0&1}.
  \]

  \medskip
  \textbf{Solución.}
  \begin{center}
  \begin{tabular}{@{}lccc@{}}
    \toprule
    \textbf{Matriz} & \textbf{Filas} & \textbf{Columnas} & \textbf{Orden / elementos} \\
    \midrule
    $A$ & 2 & 3 & $2\times3$ \quad 6 \\
    $B$ & 3 & 2 & $3\times2$ \quad 6 \\
    $C$ & 3 & 3 & $3\times3$ \quad 9 \\
    $D$ & 1 & 3 & $1\times3$ \quad 3 \\
    \bottomrule
  \end{tabular}
  \end{center}
  Observa que $A$ y $B$ tienen los mismos $6$ elementos, pero \textbf{no} son
  iguales: sus órdenes son distintos.
\end{cuadrosolucion}

\begin{cuadrosolucion}
  \small
  \textbf{Ejemplo 2.} Una matriz tiene $20$ elementos y $4$ filas. ¿Cuál es
  su orden? ¿Y si tuviera $4$ columnas?

  \medskip
  \textbf{Solución.} Como $m\cdot n=20$ y $m=4$, entonces $n=20/4=5$: el
  orden es $4\times5$. Si las columnas fueran $4$, el número de filas sería
  $20/4=5$, y el orden sería $5\times4$ (otra matriz, no la misma).
\end{cuadrosolucion}

\begin{cuadrosolucion}
  \small
  \textbf{Ejemplo 3.} ¿Cuántas matrices distintas de orden $2\times2$ se
  pueden formar usando solo los números $0$ y $1$?

  \medskip
  \textbf{Solución.} Una matriz $2\times2$ tiene $4$ posiciones y en cada una
  hay $2$ opciones independientes. Por el principio multiplicativo:
  $2^4=16$ matrices distintas (entre ellas la nula y la identidad).
\end{cuadrosolucion}

\begin{cuadroadvertencia}
  \small
  \textbf{$2\times3$ no es lo mismo que $3\times2$.} El primer número cuenta
  filas y el segundo, columnas. Una matriz $2\times3$ tiene $2$ filas y $3$
  columnas; una $3\times2$ tiene $3$ filas y $2$ columnas. Tienen el mismo
  número de elementos, pero son objetos distintos: no se pueden sumar entre
  sí y no se comportan igual al multiplicar.
\end{cuadroadvertencia}

% ============================================================
\section{Operaciones con matrices}
% ============================================================

No todas las operaciones se pueden hacer siempre: cada una tiene una
\textbf{condición de tamaño}. Esta sección estudia las cuatro operaciones
básicas en el orden en que conviene aprenderlas.

% ------------------------------------------------------------
\subsection{Multiplicación por un escalar}
% ------------------------------------------------------------

\begin{cuadroteoria}
  \small
  \textbf{Definición.} Dada una matriz $A=(a_{ij})$ y un número real $k$
  (llamado \textbf{escalar}), el producto $kA$ es la matriz que se obtiene
  multiplicando \textbf{cada} elemento de $A$ por $k$:
  \[
    kA=(k\,a_{ij}).
  \]
  El orden de la matriz \textbf{no cambia}: si $A$ es $m\times n$, entonces
  $kA$ también es $m\times n$.
\end{cuadroteoria}

\begin{cuadrosolucion}
  \small
  \textbf{Ejemplo 1.} Si $A=\mwmat{2 & -3\\ 1 & 4}$, calcular $3A$.
  \[
    3A=3\mwmat{2 & -3\\ 1 & 4}=\mwmat{3\cdot2 & 3\cdot(-3)\\ 3\cdot1 & 3\cdot4}
    =\mwmat{6 & -9\\ 3 & 12}.
  \]
\end{cuadrosolucion}

\begin{cuadrosolucion}
  \small
  \textbf{Ejemplo 2.} Si $B=\mwmat{-1 & 2\\ 0 & 5\\ 3 & -2}$, calcular $-2B$.
  \[
    -2B=\mwmat{(-2)(-1) & (-2)(2)\\ (-2)(0) & (-2)(5)\\ (-2)(3) & (-2)(-2)}
    =\mwmat{2 & -4\\ 0 & -10\\ -6 & 4}.
  \]
  El signo menos \textbf{cambia el signo de todos} los elementos.
\end{cuadrosolucion}

\begin{cuadrosolucion}
  \small
  \textbf{Ejemplo 3.} Si $C=\mwmat{\tfrac12 & -1\\[2pt] 3 & \tfrac23}$,
  calcular $6C$.
  \[
    6C=\mwmat{6\cdot\tfrac12 & 6\cdot(-1)\\[2pt] 6\cdot3 & 6\cdot\tfrac23}
    =\mwmat{3 & -6\\ 18 & 4}.
  \]
  Multiplicar por $6$ es una forma rápida de «quitar denominadores».
\end{cuadrosolucion}

\begin{cuadroteoria}
  \small
  \textbf{Propiedades.} Para matrices $A$ y $B$ del mismo orden y escalares
  $k,l\in\mathbb{R}$:
  \begin{enumerate}[leftmargin=1.4em, nosep]
    \item $k(A+B)=kA+kB$ \quad (distributiva respecto de la suma de matrices)
    \item $(k+l)A=kA+lA$ \quad (distributiva respecto de la suma de escalares)
    \item $k(lA)=(kl)A$ \quad (asociativa)
    \item $1A=A$ \qquad $0A=O$ \qquad $(-1)A=-A$
  \end{enumerate}
  Aquí $O$ es la matriz nula del mismo orden que $A$ y $-A$ es la matriz
  \textbf{opuesta}, que se obtiene cambiando el signo de todos los elementos.
\end{cuadroteoria}

\begin{cuadrosolucion}
  \small
  \textbf{Ejemplo 4.} Con $A=\mwmat{2&1\\-1&3}$ y $B=\mwmat{0&-2\\4&1}$,
  calcular $2A-3B$.

  \medskip
  \textbf{Solución.} Primero cada producto y después la resta:
  \[
    2A=\mwmat{4&2\\-2&6},\qquad
    3B=\mwmat{0&-6\\12&3},
  \]
  \[
    2A-3B=\mwmat{4-0 & 2-(-6)\\-2-12 & 6-3}
    =\mwmat{4 & 8\\ -14 & 3}.
  \]
\end{cuadrosolucion}

\begin{cuadrosolucion}
  \small
  \textbf{Ejemplo 5.} Resolver para $X$:
  $2X+\mwmat{1&2\\3&4}=\mwmat{5&6\\7&8}$.

  \medskip
  \textbf{Solución.} Se despeja como en una ecuación numérica:
  \[
    2X=\mwmat{5&6\\7&8}-\mwmat{1&2\\3&4}
      =\mwmat{4&4\\4&4}
    \quad\Rightarrow\quad
    X=\tfrac12\mwmat{4&4\\4&4}=\mwmat{2&2\\2&2}.
  \]
\end{cuadrosolucion}

\begin{cuadroadvertencia}
  \small
  \textbf{El escalar multiplica a \emph{todos}.} Un error frecuentísimo es
  multiplicar solo la primera fila o solo la diagonal. Si $kA$ se calcula
  bien, \emph{cada} entrada queda multiplicada por $k$. Tampoco se puede
  «sumar» un escalar a una matriz: $A+k$ no está definido; lo que sí existe
  es $kA$.
\end{cuadroadvertencia}

% ------------------------------------------------------------
\subsection{Suma y resta de matrices}
% ------------------------------------------------------------

\begin{cuadroteoria}
  \small
  \textbf{Condición:} $A$ y $B$ deben tener \textbf{el mismo orden}.
  \\[3pt]
  \textbf{Definición.} Si $A=(a_{ij})$ y $B=(b_{ij})$ son del mismo orden,
  \[
    A+B=(a_{ij}+b_{ij}),
    \qquad
    A-B=(a_{ij}-b_{ij}).
  \]
  Es decir, se opera \textbf{posición por posición}: cada entrada de la suma
  es la suma de las entradas que ocupan \emph{el mismo lugar}.
\end{cuadroteoria}

\begin{center}
\begin{tikzpicture}[baseline=(AA.center)]
  \matrix (AA) [matrix of math nodes, nodes={mwcelda},
                row sep=-\pgflinewidth, column sep=-\pgflinewidth] {
    a_{11} & a_{12}\\
    a_{21} & a_{22}\\
  };
  \node[right=0.3cm of AA, font=\large] {$+$};
  \matrix (BB) [matrix of math nodes, nodes={mwcelda},
                right=0.72cm of AA,
                row sep=-\pgflinewidth, column sep=-\pgflinewidth] {
    b_{11} & b_{12}\\
    b_{21} & b_{22}\\
  };
  \node[right=0.2cm of BB, font=\large] {$=$};
  \matrix (CC) [matrix of math nodes, nodes={mwcelda},
                right=0.62cm of BB,
                row sep=-\pgflinewidth, column sep=-\pgflinewidth] {
    a_{11}+b_{11} & a_{12}+b_{12}\\
    a_{21}+b_{21} & a_{22}+b_{22}\\
  };
  \draw[mwflecha] (AA-1-1.south) -- (CC-1-1.south);
  \draw[mwflecha] (BB-2-2.south) -- (CC-2-2.south);
\end{tikzpicture}
\end{center}

\begin{cuadrosolucion}
  \small
  \textbf{Ejemplo 1.} Con $A=\mwmat{1&2\\3&4}$ y $B=\mwmat{5&-1\\0&2}$:
  \[
    A+B=\mwmat{1+5 & 2+(-1)\\ 3+0 & 4+2}=\mwmat{6&1\\3&6},
  \]
  \[
    A-B=\mwmat{1-5 & 2-(-1)\\ 3-0 & 4-2}=\mwmat{-4&3\\3&2}.
  \]
\end{cuadrosolucion}

\begin{cuadrosolucion}
  \small
  \textbf{Ejemplo 2.} Con
  $A=\mwmat{2&0&-1\\3&4&5}$ y $B=\mwmat{-1&2&0\\1&-3&2}$ (ambas $2\times3$):
  \[
    A+B=\mwmat{1&2&-1\\4&1&7},
    \qquad
    A-B=\mwmat{3&-2&-1\\2&7&3}.
  \]
  Fíjate en que la resta se hace \emph{casilla por casilla}, conservando la
  posición.
\end{cuadrosolucion}

\begin{cuadrosolucion}
  \small
  \textbf{Ejemplo 3.} Determinar $a$ y $b$ para que
  \[
    \mwmat{a&2\\3&b}+\mwmat{1&-2\\-3&4}=\mwmat{4&0\\0&7}.
  \]
  \textbf{Solución.} Sumando posición por posición:
  \[
    a+1=4\Rightarrow a=3,\qquad 3+(-3)=0\ \checkmark,\qquad b+4=7\Rightarrow b=3.
  \]
  Además $2+(-2)=0$, que coincide con la entrada correspondiente del lado
  derecho: la solución es consistente.
\end{cuadrosolucion}

\begin{cuadrosolucion}
  \small
  \textbf{Ejemplo 4.} Con $A=\mwmat{1&0\\2&-1}$,
  $B=\mwmat{3&1\\-1&2}$ y $C=\mwmat{0&4\\1&0}$, calcular $2A-B+C$.

  \medskip
  \textbf{Solución.}
  \[
    2A=\mwmat{2&0\\4&-2},\qquad
    2A-B=\mwmat{2-3 & 0-1\\ 4-(-1) & -2-2}=\mwmat{-1&-1\\5&-4},
  \]
  \[
    2A-B+C=\mwmat{-1+0 & -1+4\\ 5+1 & -4+0}=\mwmat{-1&3\\6&-4}.
  \]
\end{cuadrosolucion}

\begin{cuadroteoria}
  \small
  \textbf{Propiedades.} Para matrices del mismo orden:
  \begin{enumerate}[leftmargin=1.4em, nosep]
    \item $A+B=B+A$ \quad (conmutativa)
    \item $(A+B)+C=A+(B+C)$ \quad (asociativa)
    \item $A+O=A$ \quad ($O$ es el neutro aditivo)
    \item $A+(-A)=O$ \quad ($-A$ es el opuesto de $A$)
    \item $A-B=A+(-B)$
  \end{enumerate}
  La resta \textbf{no} es conmutativa: en general $A-B\neq B-A$ (de hecho,
  $B-A=-(A-B)$).
\end{cuadroteoria}

\begin{cuadroadvertencia}
  \small
  \textbf{Dos errores clásicos:}
  \begin{itemize}[leftmargin=1.2em, nosep]
    \item \textbf{Sumar matrices de distinto orden.} Si $A$ es $2\times3$ y
          $B$ es $3\times2$, $A+B$ \textbf{no existe}. No importa que ambas
          tengan $6$ elementos: la suma se define posición por posición y las
          posiciones no coinciden.
    \item \textbf{Olvidar los signos.} Restar una matriz equivale a sumar su
          opuesta: cada entrada de $B$ cambia de signo. El error típico es
          cambiar el signo solo en la primera fila.
  \end{itemize}
\end{cuadroadvertencia}

% ------------------------------------------------------------
\subsection{Sistemas de ecuaciones con matrices}
% ------------------------------------------------------------

Empezamos por el caso base ---\textbf{dos ecuaciones con dos incógnitas}--- y
al final de la sección lo extendemos a \textbf{tres ecuaciones con tres
incógnitas}.

Un sistema de dos ecuaciones lineales con dos incógnitas,
\[
  \begin{cases}
    a_{11}x+a_{12}y=c_1,\\
    a_{21}x+a_{22}y=c_2,
  \end{cases}
\]
se escribe de forma compacta como una \textbf{ecuación matricial}.

\begin{cuadroteoria}
  \small
  \textbf{Forma matricial.} El sistema anterior equivale a
  \[
    \underbrace{\mwmat{a_{11}&a_{12}\\a_{21}&a_{22}}}_{A}
    \underbrace{\mwmat{x\\y}}_{\mathbf{x}}
    =
    \underbrace{\mwmat{c_1\\c_2}}_{\mathbf{b}},
    \qquad\text{es decir,}\qquad A\mathbf{x}=\mathbf{b}.
  \]
  \begin{itemize}[leftmargin=1.2em, nosep]
    \item $A$ es la \textbf{matriz de coeficientes} (orden $2\times2$);
    \item $\mathbf{x}$ es el \textbf{vector de incógnitas} ($2\times1$);
    \item $\mathbf{b}$ es el \textbf{vector de constantes} ($2\times1$).
  \end{itemize}
\end{cuadroteoria}

Cada \textbf{fila} de $A$ contiene los coeficientes de una ecuación; cada
\textbf{columna}, los coeficientes de una misma incógnita. Al añadir la
columna de las constantes se obtiene la \textbf{matriz aumentada}:
\[
  \left[A\mid\mathbf{b}\right]
  =
  \mwaug{cc|c}{a_{11}&a_{12}&c_1\\ a_{21}&a_{22}&c_2}.
\]

\begin{cuadrosolucion}
  \small
  \textbf{Ejemplo 1.} Escribir en forma matricial y en matriz aumentada el
  sistema
  \[
    \begin{cases}
      2x+3y=8,\\
      x+4y=9.
    \end{cases}
  \]

  \medskip
  \textbf{Solución.}
  \[
    \mwmat{2&3\\1&4}\mwmat{x\\y}=\mwmat{8\\9},
    \qquad
    \mwaug{cc|c}{2&3&8\\1&4&9}.
  \]
  \textbf{Verificación de la solución} (que obtendremos más adelante):
  $x=1,\;y=2$ da $2(1)+3(2)=8$ y $1+4(2)=9$. \checkmark
\end{cuadrosolucion}

\begin{cuadrosolucion}
  \small
  \textbf{Ejemplo 2.} Escribir el sistema asociado a
  $\mwaug{cc|c}{3&-1&5\\1&2&4}$ y comprobar si $(2,1)$ es solución.

  \medskip
  \textbf{Solución.} Las filas dan las ecuaciones
  \[
    \begin{cases}
      3x-y=5,\\
      x+2y=4.
    \end{cases}
  \]
  Con $x=2,\;y=1$: $3(2)-1=5$ \checkmark\ y $2+2(1)=4$ \checkmark. Sí es
  solución (de hecho es la única).
\end{cuadrosolucion}

\subsubsection*{El mismo esquema en $3\times3$}

Para tres ecuaciones con tres incógnitas la construcción es idéntica:
\[
  \begin{cases}
    a_{11}x+a_{12}y+a_{13}z=c_1,\\
    a_{21}x+a_{22}y+a_{23}z=c_2,\\
    a_{31}x+a_{32}y+a_{33}z=c_3,
  \end{cases}
  \quad\Longleftrightarrow\quad
  \underbrace{\mwmat{a_{11}&a_{12}&a_{13}\\a_{21}&a_{22}&a_{23}\\a_{31}&a_{32}&a_{33}}}_{A}
  \underbrace{\mwmat{x\\y\\z}}_{\mathbf{x}}
  =
  \underbrace{\mwmat{c_1\\c_2\\c_3}}_{\mathbf{b}}.
\]
La matriz de coeficientes es $3\times3$, los vectores son $3\times1$ y la
matriz aumentada resulta de orden $3\times4$:
\[
  \left[A\mid\mathbf{b}\right]=
  \mwaug{ccc|c}{a_{11}&a_{12}&a_{13}&c_1\\
                a_{21}&a_{22}&a_{23}&c_2\\
                a_{31}&a_{32}&a_{33}&c_3}.
\]

\begin{cuadrosolucion}
  \small
  \textbf{Ejemplo 3.} Escribir en forma matricial y en matriz aumentada
  \[
    \begin{cases}
      x+y+z=6,\\
      2x-y+z=3,\\
      x+2y-z=2.
    \end{cases}
  \]

  \medskip
  \textbf{Solución.}
  \[
    \mwmat{1&1&1\\2&-1&1\\1&2&-1}\mwmat{x\\y\\z}=\mwmat{6\\3\\2},
    \qquad
    \mwaug{ccc|c}{1&1&1&6\\2&-1&1&3\\1&2&-1&2}.
  \]
  Este sistema tiene solución única: $(x,y,z)=(1,2,3)$. Compruébalo:
  $1+2+3=6$, $2-2+3=3$ y $1+4-3=2$. Lo resolveremos por Gauss en la §6.
\end{cuadrosolucion}

\begin{cuadroadvertencia}
  \small
  Al pasar a $3\times3$ aparecen dos errores nuevos: \textbf{olvidar una
  incógnita} en alguna ecuación (su coeficiente es $0$, pero la columna
  existe) y \textbf{desordenar las incógnitas} entre ecuaciones. Antes de
  escribir la matriz, alinea mentalmente las tres ecuaciones: $x$, luego $y$,
  luego $z$, y al final la constante.
\end{cuadroadvertencia}

\subsubsection*{Casos especiales: $2\times3$ y $3\times2$}

El número de incógnitas no siempre coincide con el número de ecuaciones.
Cuando eso ocurre aparecen dos situaciones que conviene reconocer.

\begin{cuadroteoria}
  \small
  \textbf{Caso $2\times3$ (dos ecuaciones, tres incógnitas).} La matriz de
  coeficientes es $2\times3$; la aumentada, $2\times4$. Normalmente el
  sistema tiene \textbf{infinitas soluciones}: una incógnita queda libre y
  las demás se expresan en función de ella.
  \\[4pt]
  \textbf{Ejemplo:}
  \[
    \begin{cases}
      x+y+z=3,\\
      2x-y+z=0.
    \end{cases}
    \qquad
    \mwaug{ccc|c}{1&1&1&3\\2&-1&1&0}.
  \]
  Tomando $z$ como parámetro: sumando las dos ecuaciones se obtiene
  $3x+2z=3$, es decir
  \[
    x=1-\tfrac{2}{3}z,
    \qquad
    y=2-\tfrac{1}{3}z,
    \qquad z\in\mathbb{R}.
  \]
  Por ejemplo, con $z=0$: $(x,y,z)=(1,2,0)$; con $z=3$: $(-1,1,3)$.
\end{cuadroteoria}

\begin{cuadroteoria}
  \small
  \textbf{Caso $3\times2$ (tres ecuaciones, dos incógnitas).} La matriz de
  coeficientes es $3\times2$; la aumentada, $3\times3$. Hay \textbf{más
  ecuaciones que incógnitas}, así que el sistema suele ser
  \textbf{incompatible}: la tercera ecuación aporta una condición extra que
  las dos primeras no cumplen.
  \\[4pt]
  \textbf{Ejemplo compatible:}
  \[
    \begin{cases}
      x+y=2,\\
      x-y=0,\\
      2x+y=3.
    \end{cases}
  \]
  De las dos primeras, $x=1$, $y=1$. La tercera se cumple: $2(1)+1=3$.
  Solución única: $(1,1)$.
  \\[4pt]
  \textbf{El mismo sistema con otra tercera ecuación,}
  $2x+2y=5$, sería \textbf{incompatible}: con $x=1$, $y=1$ se obtiene
  $4\neq5$.
\end{cuadroteoria}

\begin{cuadroadvertencia}
  \small
  \textbf{Antes de escribir la matriz aumentada:}
  \begin{enumerate}[leftmargin=1.4em, nosep]
    \item ordena las ecuaciones y las incógnitas siempre en el mismo orden
          (por ejemplo, $x$ antes que $y$);
    \item si una incógnita no aparece en una ecuación, su coeficiente es
          $0$ (no se omite la columna);
    \item si una ecuación está escrita como $3x=7-2y$, pásala primero a la
          forma $3x+2y=7$;
    \item la columna de las constantes va \textbf{siempre al final}, separada
          por una barra.
  \end{enumerate}
\end{cuadroadvertencia}

% ------------------------------------------------------------
\subsection{Multiplicación de matrices}
% ------------------------------------------------------------

\begin{cuadroteoria}
  \small
  \textbf{Condición de conformabilidad.} El producto $AB$ está definido
  únicamente si
  \[
    \text{número de columnas de }A
    =
    \text{número de filas de }B.
  \]
  Si $A$ es $m\times n$ y $B$ es $n\times p$, el producto $AB$ es de orden
  $m\times p$:
  \[
    (m\times n)\cdot(n\times p)\;\longrightarrow\;(m\times p).
  \]
\end{cuadroteoria}

\begin{cuadroteoria}
  \small
  \textbf{Definición.} Si $A=(a_{ij})$ es $m\times n$ y $B=(b_{ij})$ es
  $n\times p$, entonces $C=AB=(c_{ij})$ es $m\times p$ con
  \[
    c_{ij}=\sum_{k=1}^{n}a_{ik}b_{kj}
    =a_{i1}b_{1j}+a_{i2}b_{2j}+\cdots+a_{in}b_{nj}.
  \]
  En palabras: para obtener la entrada $c_{ij}$ se toma la \textbf{fila $i$
  de $A$} y la \textbf{columna $j$ de $B$}, se multiplican sus elementos
  \emph{uno a uno en el mismo orden} y se suman los productos.
\end{cuadroteoria}

\begin{center}
\begin{tikzpicture}[baseline=(A.center)]
  \matrix (A) [matrix of math nodes, nodes={mwcelda},
               row sep=-\pgflinewidth, column sep=-\pgflinewidth] {
    a_{11} & a_{12} & a_{13}\\
    a_{21} & a_{22} & a_{23}\\
  };
  \matrix (B) [matrix of math nodes, nodes={mwcelda},
               right=1.1cm of A,
               row sep=-\pgflinewidth, column sep=-\pgflinewidth] {
    b_{11} & b_{12}\\
    b_{21} & b_{22}\\
    b_{31} & b_{32}\\
  };
  \matrix (C) [matrix of math nodes, nodes={mwcelda},
               right=1.1cm of B,
               row sep=-\pgflinewidth, column sep=-\pgflinewidth] {
    c_{11} & c_{12}\\
    c_{21} & c_{22}\\
  };
  \node at ($(A.east)!0.5!(B.west)$) {$\cdot$};
  \node at ($(B.east)!0.5!(C.west)$) {$=$};
  \draw[mwnaranja, thick, rounded corners=2pt]
    (A-2-1.north west) rectangle (A-2-3.south east);
  \draw[mwmagenta, thick, rounded corners=2pt]
    (B-1-1.north west) rectangle (B-3-1.south east);
  \draw[mwvioleta, thick, rounded corners=2pt]
    (C-2-1.north west) rectangle (C-2-1.south east);
  \node[below=0.3cm of A, font=\scriptsize, mwnaranja] {$2\times3$};
  \node[below=0.3cm of B, font=\scriptsize, mwmagenta] {$3\times2$};
  \node[below=0.3cm of C, font=\scriptsize, mwvioleta] {$2\times2$};
\end{tikzpicture}
\end{center}

\begin{center}
\small
$c_{21}=a_{21}b_{11}+a_{22}b_{21}+a_{23}b_{31}$
\qquad (fila $2$ de $A$ por columna $1$ de $B$)
\end{center}

\begin{cuadrosolucion}
  \small
  \textbf{Ejemplo 1 ($2\times2$ por $2\times2$).} Con
  $A=\mwmat{1&2\\3&4}$ y $B=\mwmat{5&6\\7&8}$:
  \[
    AB=\mwmat{1\cdot5+2\cdot7 & 1\cdot6+2\cdot8\\
              3\cdot5+4\cdot7 & 3\cdot6+4\cdot8}
      =\mwmat{19&22\\43&50}.
  \]
\end{cuadrosolucion}

\begin{cuadrosolucion}
  \small
  \textbf{Ejemplo 2 ($2\times3$ por $3\times2$).} Con
  $A=\mwmat{1&2&3\\4&5&6}$ y $B=\mwmat{1&0\\0&1\\1&1}$:
  \[
    AB=\mwmat{
      1\cdot1+2\cdot0+3\cdot1 & 1\cdot0+2\cdot1+3\cdot1\\
      4\cdot1+5\cdot0+6\cdot1 & 4\cdot0+5\cdot1+6\cdot1}
      =\mwmat{4&5\\10&11}.
  \]
  El resultado es $2\times2$: se «cancelan» las dimensiones internas.
\end{cuadrosolucion}

\begin{cuadrosolucion}
  \small
  \textbf{Ejemplo 3 ($3\times2$ por $2\times3$).} Con
  $A=\mwmat{1&2\\3&4\\5&6}$ y $B=\mwmat{1&0&1\\0&1&1}$:
  \[
    AB=\mwmat{
      1\cdot1+2\cdot0 & 1\cdot0+2\cdot1 & 1\cdot1+2\cdot1\\
      3\cdot1+4\cdot0 & 3\cdot0+4\cdot1 & 3\cdot1+4\cdot1\\
      5\cdot1+6\cdot0 & 5\cdot0+6\cdot1 & 5\cdot1+6\cdot1}
      =\mwmat{1&2&3\\3&4&7\\5&6&11}.
  \]
  Ahora el resultado es $3\times3$.
\end{cuadrosolucion}

\begin{cuadrosolucion}
  \small
  \textbf{Ejemplo 4 (el orden importa).} Con
  $A=\mwmat{1&2\\3&4}$ y $B=\mwmat{0&1\\1&0}$:
  \[
    AB=\mwmat{2&1\\4&3},
    \qquad
    BA=\mwmat{3&4\\1&2}.
  \]
  Por lo tanto $AB\neq BA$. La multiplicación de matrices \textbf{no es
  conmutativa}.
\end{cuadrosolucion}

\begin{cuadroteoria}
  \small
  \textbf{Propiedades.} Siempre que las operaciones estén definidas:
  \begin{enumerate}[leftmargin=1.4em, nosep]
    \item $(AB)C=A(BC)$ \quad (asociativa)
    \item $A(B+C)=AB+AC$ \quad y \quad $(A+B)C=AC+BC$ \quad (distributivas)
    \item $AI=IA=A$ \quad ($I$ es la identidad del orden adecuado)
    \item $AO=OA=O$
    \item $(kA)B=k(AB)=A(kB)$
    \item \textbf{No} es conmutativa: en general $AB\neq BA$.
    \item \textbf{No} se puede «cancelar»: $AB=AC$ \emph{no} implica $B=C$.
    \item $AB=O$ \emph{no} implica $A=O$ ni $B=O$.
  \end{enumerate}
\end{cuadroteoria}

\begin{cuadrosolucion}
  \small
  \textbf{Ejemplo 5 ($AB=O$ con factores no nulos).} Con
  $A=\mwmat{1&1\\2&2}$ y $B=\mwmat{1&-1\\-1&1}$:
  \[
    AB=\mwmat{1-1 & -1+1\\ 2-2 & -2+2}=\mwmat{0&0\\0&0}=O,
  \]
  y sin embargo $A\neq O$ y $B\neq O$. Moraleja: en matrices, un producto
  nulo \textbf{no} permite concluir que algún factor sea nulo.
\end{cuadrosolucion}

\begin{cuadrosolucion}
  \small
  \textbf{Ejemplo 6 (potencias).} Si $A=\mwmat{1&1\\0&1}$, calcular $A^2$ y
  $A^3$, y conjeturar $A^n$.
  \[
    A^2=\mwmat{1&1\\0&1}\mwmat{1&1\\0&1}=\mwmat{1&2\\0&1},
    \qquad
    A^3=A^2A=\mwmat{1&3\\0&1}.
  \]
  En general,
  \[
    A^n=\mwmat{1&n\\0&1}\qquad(n\ge1).
  \]
\end{cuadrosolucion}

\begin{cuadroadvertencia}
  \small
  \textbf{La multiplicación no se hace «entrada con entrada».} El producto
  $AB$ \textbf{no} es la matriz cuyas entradas son $a_{ij}b_{ij}$. Hay que
  combinar \emph{toda} una fila de $A$ con \emph{toda} una columna de $B$.
  Además, antes de empezar, \textbf{verifica la condición de tamaño}: si
  escribes el producto sin comprobarla, lo más probable es que el cálculo sea
  inválido desde el principio.
\end{cuadroadvertencia}

% ============================================================
\section{Matrices equivalentes}
% ============================================================

\begin{cuadroteoria}
  \small
  \textbf{Definición.} Dos matrices $A$ y $B$ del mismo orden son
  \textbf{equivalentes} (se escribe $A\sim B$) si $B$ se puede obtener a
  partir de $A$ mediante una cantidad finita de \textbf{operaciones
  elementales de fila}:
  \begin{enumerate}[leftmargin=1.4em, nosep]
    \item \textbf{Intercambiar} dos filas: $F_i\leftrightarrow F_j$.
    \item \textbf{Multiplicar} una fila por un escalar \emph{no nulo}:
          $F_i\leftarrow kF_i$, con $k\neq0$.
    \item \textbf{Sumar} a una fila un múltiplo de otra:
          $F_i\leftarrow F_i+kF_j$ (con $i\neq j$).
  \end{enumerate}
  La equivalencia se aplica \textbf{sobre las filas}, nunca sobre las
  columnas.
  \\[3pt]
  \textbf{Notación.} $F_i$ denota la \textbf{fila} $i$ (inicial de
  \emph{fila}). Así, $F_1\leftrightarrow F_2$ se lee «intercambiar la fila
  $1$ con la fila $2$», y $F_2\leftarrow F_2-3F_1$ se lee «a la fila $2$ se
  le resta el triple de la fila $1$».
\end{cuadroteoria}

No hay que confundir tres ideas parecidas:

\begin{center}
\small
\begin{tabular}{@{}p{0.24\linewidth} p{0.66\linewidth}@{}}
  \toprule
  \textbf{Concepto} & \textbf{Significado} \\
  \midrule
  Matrices \textbf{iguales} & Todos los elementos coinciden posición por
    posición. Es una igualdad estricta. \\
  Matrices \textbf{equivalentes} & Se pasa de una a otra con operaciones
    elementales de fila. Pueden tener números distintos. \\
  \textbf{Sistemas equivalentes} & Tienen exactamente el mismo conjunto
    solución, aunque sus ecuaciones se vean diferentes. \\
  \bottomrule
\end{tabular}
\end{center}

\begin{cuadrosolucion}
  \small
  \textbf{Ejemplo 1.} Mostrar que
  $A=\mwmat{1&2\\2&5}$ es equivalente a $I_2$.

  \medskip
  \textbf{Solución.} Encadenamos operaciones elementales:
  \[
    \mwmat{1&2\\2&5}
    \;\xrightarrow{\;F_2\leftarrow F_2-2F_1\;}
    \mwmat{1&2\\0&1}
    \;\xrightarrow{\;F_1\leftarrow F_1-2F_2\;}
    \mwmat{1&0\\0&1}=I_2.
  \]
  Por lo tanto $A\sim I_2$.
\end{cuadrosolucion}

\begin{cuadrosolucion}
  \small
  \textbf{Ejemplo 2 (equivalencia y sistemas).} Resolver
  \[
    \begin{cases}
      x+2y=8,\\
      3x-y=3,
    \end{cases}
  \]
  trabajando solo con la matriz aumentada.

  \medskip
  \textbf{Solución.} Cada operación de fila produce una matriz aumentada
  \emph{equivalente}, es decir, un sistema con las mismas soluciones:
  \[
    \mwaug{cc|c}{1&2&8\\3&-1&3}
    \xrightarrow{\;F_2\leftarrow F_2-3F_1\;}
    \mwaug{cc|c}{1&2&8\\0&-7&-21}
    \xrightarrow{\;F_2\leftarrow -\frac17F_2\;}
    \mwaug{cc|c}{1&2&8\\0&1&3}.
  \]
  La última matriz representa el sistema
  \[
    \begin{cases}
      x+2y=8,\\
      y=3,
    \end{cases}
    \quad\Rightarrow\quad
    x=8-2(3)=2.
  \]
  Solución: $(x,y)=(2,3)$.
\end{cuadrosolucion}

\begin{center}
\begin{tikzpicture}[baseline=(M1.center)]
  \matrix (M1) [matrix of math nodes, nodes={mwcelda},
                row sep=-\pgflinewidth, column sep=-\pgflinewidth] {
    1 & 2 & 8\\
    3 & -1 & 3\\
  };
  \matrix (M2) [matrix of math nodes, nodes={mwcelda},
                right=2.4cm of M1,
                row sep=-\pgflinewidth, column sep=-\pgflinewidth] {
    1 & 2 & 8\\
    0 & -7 & -21\\
  };
  \matrix (M3) [matrix of math nodes, nodes={mwcelda},
                right=2.4cm of M2,
                row sep=-\pgflinewidth, column sep=-\pgflinewidth] {
    1 & 2 & 8\\
    0 & 1 & 3\\
  };
  \draw[mwflecha] (M1) -- node[above, font=\scriptsize, mwnaranja]
    {$F_2\leftarrow F_2-3F_1$} (M2);
  \draw[mwflecha] (M2) -- node[above, font=\scriptsize, mwnaranja]
    {$F_2\leftarrow -\frac17F_2$} (M3);
  \node[below=0.25cm of M1, font=\scriptsize, mwgris] {matriz original};
  \node[below=0.25cm of M3, font=\scriptsize, mwgris] {forma escalonada};
\end{tikzpicture}
\end{center}

\begin{cuadroteoria}
  \small
  \textbf{Propiedad fundamental.} Si dos matrices \textbf{aumentadas}
  $[A\mid\mathbf{b}]$ y $[A'\mid\mathbf{b}']$ son equivalentes por filas,
  entonces los sistemas $A\mathbf{x}=\mathbf{b}$ y
  $A'\mathbf{x}=\mathbf{b}'$ tienen \textbf{el mismo conjunto solución}.
  Esta es la razón por la que el método de Gauss funciona: no cambiamos el
  sistema, solo su «disfraz».
\end{cuadroteoria}

\begin{cuadroadvertencia}
  \small
  \textbf{Lo que sí y lo que no conserva la equivalencia:}
  \begin{itemize}[leftmargin=1.2em, nosep]
    \item \textbf{Sí} se pueden intercambiar filas, multiplicar una fila por
          un número distinto de cero y sumar múltiplos de filas.
    \item \textbf{No} se puede multiplicar una fila por $0$: se destruye
          información y el sistema deja de ser equivalente.
    \item \textbf{No} se opera sobre columnas: intercambiar o combinar
          columnas cambia el significado de las incógnitas y puede alterar
          el conjunto solución.
  \end{itemize}
\end{cuadroadvertencia}

% ============================================================
\section{Operaciones que producen matrices equivalentes}
% ============================================================

Las tres operaciones elementales de fila son las únicas herramientas
permitidas para transformar una matriz sin cambiar el sistema que representa.
Veamos cada una con detalle.

\subsection*{Operación 1: intercambiar dos filas}

\[
  \mwmat{1&2\\3&4}
  \;\xrightarrow{\;F_1\leftrightarrow F_2\;}
  \mwmat{3&4\\1&2}.
\]
Se usa cuando el pivote de la fila en curso es $0$ y conviene «subir» otra
fila. El determinante cambia de signo, pero el conjunto solución del sistema
no cambia.

\subsection*{Operación 2: multiplicar una fila por un escalar no nulo}

\[
  \mwmat{2&4&10\\1&3&7}
  \;\xrightarrow{\;F_1\leftarrow \frac12F_1\;}
  \mwmat{1&2&5\\1&3&7}.
\]
Se usa para convertir el pivote en $1$ o para simplificar denominadores. El
escalar \textbf{debe ser distinto de cero}: multiplicar por $0$ convierte la
ecuación en $0=0$ y pierde información.

\subsection*{Operación 3: sumar a una fila un múltiplo de otra}

\[
  \mwmat{1&2&5\\2&5&12}
  \;\xrightarrow{\;F_2\leftarrow F_2-2F_1\;}
  \mwmat{1&2&5\\0&1&2}.
\]
Es la operación que \textbf{elimina} coeficientes y produce los ceros de la
forma escalonada. También se permite en sentido inverso, sumando un múltiplo
de una fila \emph{superior} a una inferior o al revés.

\begin{cuadroteoria}
  \small
  \textbf{Forma escalonada (por filas).} Una matriz está en forma
  escalonada si:
  \begin{enumerate}[leftmargin=1.4em, nosep]
    \item todas las filas nulas, si existen, están al final;
    \item el primer elemento no nulo de cada fila (el \textbf{pivote}) está
          a la derecha del pivote de la fila anterior;
    \item debajo de cada pivote todos los elementos son $0$.
  \end{enumerate}
  Está en \textbf{forma escalonada reducida} si además cada pivote vale $1$
  y es el único elemento no nulo de su columna.
\end{cuadroteoria}

\begin{cuadrosolucion}
  \small
  \textbf{Ejemplo 1.} Reducir
  $\mwmat{2&4&10\\1&3&7}$ a forma escalonada y luego a forma escalonada
  reducida.

  \medskip
  \textbf{Solución.}
  \[
    \mwmat{2&4&10\\1&3&7}
    \xrightarrow{F_1\leftrightarrow F_2}
    \mwmat{1&3&7\\2&4&10}
    \xrightarrow{F_2\leftarrow F_2-2F_1}
    \mwmat{1&3&7\\0&-2&-4}
  \]
  \[
    \xrightarrow{F_2\leftarrow -\frac12F_2}
    \mwmat{1&3&7\\0&1&2}
    \xrightarrow{F_1\leftarrow F_1-3F_2}
    \mwmat{1&0&1\\0&1&2}.
  \]
  La penúltima matriz está en forma escalonada; la última, en forma
  escalonada reducida.
\end{cuadrosolucion}

\begin{cuadrosolucion}
  \small
  \textbf{Ejemplo 2 (combinación de las tres operaciones).} Reducir
  \[
    \mwmat{0&2&4\\3&6&9}.
  \]
  \textbf{Solución.} Primero se necesita un pivote no nulo en la primera
  fila, así que se intercambian las filas:
  \[
    \mwmat{0&2&4\\3&6&9}
    \xrightarrow{F_1\leftrightarrow F_2}
    \mwmat{3&6&9\\0&2&4}
    \xrightarrow{F_1\leftarrow \frac13F_1}
    \mwmat{1&2&3\\0&2&4}
    \xrightarrow{F_2\leftarrow \frac12F_2}
    \mwmat{1&2&3\\0&1&2}.
  \]
\end{cuadrosolucion}

\begin{cuadroteoria}
  \small
  \textbf{La equivalencia es una relación de equivalencia.} Para matrices
  del mismo orden:
  \begin{enumerate}[leftmargin=1.4em, nosep]
    \item \textbf{Reflexiva:} $A\sim A$ (no hacer nada es una secuencia
          válida de operaciones).
    \item \textbf{Simétrica:} si $A\sim B$, entonces $B\sim A$ (cada
          operación elemental se puede deshacer con otra operación
          elemental).
    \item \textbf{Transitiva:} si $A\sim B$ y $B\sim C$, entonces
          $A\sim C$ (se concatenan las secuencias de operaciones).
  \end{enumerate}
\end{cuadroteoria}

\begin{cuadropasos}
  \small
  \textbf{Protocolo para reducir sin perder información:}
  \begin{enumerate}[leftmargin=1.4em, nosep]
    \item Trabaja \textbf{de izquierda a derecha} y de arriba hacia abajo.
    \item En la columna actual, consigue un pivote no nulo (si es necesario,
          intercambia filas).
    \item Convierte el pivote en $1$ dividiendo toda la fila.
    \item Anula los elementos \textbf{debajo} del pivote sumando múltiplos
          adecuados de su fila.
    \item Pasa a la siguiente columna y repite.
    \item Si buscas la forma reducida, vuelve hacia arriba anulando también
          los elementos \textbf{encima} de cada pivote.
    \item Anota cada operación junto a la flecha: te permite revisar y
          detectar errores.
  \end{enumerate}
\end{cuadropasos}

\begin{cuadroadvertencia}
  \small
  \textbf{Errores que rompen la equivalencia:}
  \begin{itemize}[leftmargin=1.2em, nosep]
    \item multiplicar una fila por $0$ (se pierde una ecuación);
    \item sumar a una fila \emph{ella misma} ($F_i\leftarrow F_i+F_i$ no
          aporta nada nuevo y puede confundir);
    \item operar con columnas como si fueran filas;
    \item cambiar un signo «solo en una parte» de la fila: la operación se
          aplica a \textbf{toda} la fila, incluida la columna de constantes.
  \end{itemize}
\end{cuadroadvertencia}

% ============================================================
\section{Sistemas de ecuaciones por Gauss y Gauss-Jordan}
% ============================================================

Ambos métodos trabajan sobre la \textbf{matriz aumentada} y usan únicamente
operaciones elementales de fila. La diferencia está en \emph{hasta dónde}
se reduce la matriz.

\begin{cuadroteoria}
  \small
  \textbf{Método de Gauss (eliminación gaussiana).}
  \begin{enumerate}[leftmargin=1.4em, nosep]
    \item Escribir el sistema como matriz aumentada.
    \item Aplicar operaciones de fila para llevarla a \textbf{forma
          escalonada} (ceros debajo de cada pivote).
    \item Resolver por \textbf{sustitución hacia atrás}: de la última
          ecuación hacia la primera.
  \end{enumerate}
  \textbf{Método de Gauss-Jordan.}
  \begin{enumerate}[leftmargin=1.4em, nosep]
    \item Llevar la matriz aumentada hasta \textbf{forma escalonada
          reducida} (pivotes iguales a $1$ y ceros también encima).
    \item Leer la solución \textbf{directamente}, sin sustituir.
  \end{enumerate}
\end{cuadroteoria}

\begin{cuadrosolucion}
  \small
  \textbf{Ejemplo 1 (solución única, por Gauss).} Resolver
  \[
    \begin{cases}
      2x+y=5,\\
      x-y=1.
    \end{cases}
  \]

  \medskip
  \textbf{Solución.}
  \[
    \mwaug{cc|c}{2&1&5\\1&-1&1}
    \xrightarrow{F_1\leftrightarrow F_2}
    \mwaug{cc|c}{1&-1&1\\2&1&5}
    \xrightarrow{F_2\leftarrow F_2-2F_1}
    \mwaug{cc|c}{1&-1&1\\0&3&3}
  \]
  \[
    \xrightarrow{F_2\leftarrow \frac13F_2}
    \mwaug{cc|c}{1&-1&1\\0&1&1}.
  \]
  Sustitución hacia atrás: $y=1$; luego $x-1=1$, es decir $x=2$.
  \[
    (x,y)=(2,1).
  \]
  \textbf{Comprobación:} $2(2)+1=5$ \checkmark\ y $2-1=1$ \checkmark.
\end{cuadrosolucion}

\begin{cuadrosolucion}
  \small
  \textbf{Ejemplo 2 (sistema incompatible).} Resolver
  \[
    \begin{cases}
      x+y=2,\\
      2x+2y=5.
    \end{cases}
  \]

  \medskip
  \textbf{Solución.}
  \[
    \mwaug{cc|c}{1&1&2\\2&2&5}
    \xrightarrow{F_2\leftarrow F_2-2F_1}
    \mwaug{cc|c}{1&1&2\\0&0&1}.
  \]
  La segunda fila representa la ecuación $0x+0y=1$, es decir $0=1$, que es
  falsa. El sistema \textbf{no tiene solución} (es incompatible).
\end{cuadrosolucion}

\begin{cuadrosolucion}
  \small
  \textbf{Ejemplo 3 (infinitas soluciones).} Resolver
  \[
    \begin{cases}
      x+2y=4,\\
      2x+4y=8.
    \end{cases}
  \]

  \medskip
  \textbf{Solución.}
  \[
    \mwaug{cc|c}{1&2&4\\2&4&8}
    \xrightarrow{F_2\leftarrow F_2-2F_1}
    \mwaug{cc|c}{1&2&4\\0&0&0}.
  \]
  La segunda ecuación desaparece: solo queda $x+2y=4$. La incógnita $y$
  queda libre:
  \[
    y=t,\qquad x=4-2t,\qquad t\in\mathbb{R}.
  \]
  Hay \textbf{infinitas soluciones}: $(x,y)=(4-2t,\,t)$.
\end{cuadrosolucion}

\begin{cuadrosolucion}
  \small
  \textbf{Ejemplo 4 (Gauss-Jordan).} Resolver
  \[
    \begin{cases}
      3x+2y=12,\\
      x-y=-1.
    \end{cases}
  \]

  \medskip
  \textbf{Solución.}
  \[
    \mwaug{cc|c}{3&2&12\\1&-1&-1}
    \xrightarrow{F_1\leftrightarrow F_2}
    \mwaug{cc|c}{1&-1&-1\\3&2&12}
    \xrightarrow{F_2\leftarrow F_2-3F_1}
    \mwaug{cc|c}{1&-1&-1\\0&5&15}
  \]
  \[
    \xrightarrow{F_2\leftarrow \frac15F_2}
    \mwaug{cc|c}{1&-1&-1\\0&1&3}
    \xrightarrow{F_1\leftarrow F_1+F_2}
    \mwaug{cc|c}{1&0&2\\0&1&3}.
  \]
  La última matriz es la identidad a la izquierda, así que la solución se
  lee directamente:
  \[
    (x,y)=(2,3).
  \]
\end{cuadrosolucion}

\begin{cuadrosolucion}
  \small
  \textbf{Ejemplo 5 (Gauss-Jordan paso a paso).} Resolver
  \[
    \begin{cases}
      2x+3y=7,\\
      x+2y=4.
    \end{cases}
  \]

  \medskip
  \textbf{Solución.}
  \[
    \mwaug{cc|c}{2&3&7\\1&2&4}
    \xrightarrow{F_1\leftrightarrow F_2}
    \mwaug{cc|c}{1&2&4\\2&3&7}
    \xrightarrow{F_2\leftarrow F_2-2F_1}
    \mwaug{cc|c}{1&2&4\\0&-1&-1}
  \]
  \[
    \xrightarrow{F_2\leftarrow -F_2}
    \mwaug{cc|c}{1&2&4\\0&1&1}
    \xrightarrow{F_1\leftarrow F_1-2F_2}
    \mwaug{cc|c}{1&0&2\\0&1&1}.
  \]
  Por lo tanto $(x,y)=(2,1)$.
\end{cuadrosolucion}

\begin{cuadrosolucion}
  \small
  \textbf{Ejemplo 6 (sistema $3\times3$ por Gauss).} Resolver
  \[
    \begin{cases}
      x+y+z=6,\\
      2x-y+z=3,\\
      x+2y-z=2.
    \end{cases}
  \]

  \medskip
  \textbf{Solución.} Se anula la primera columna por debajo del pivote:
  \[
    \mwaug{ccc|c}{1&1&1&6\\2&-1&1&3\\1&2&-1&2}
    \xrightarrow{F_2\leftarrow F_2-2F_1}
    \mwaug{ccc|c}{1&1&1&6\\0&-3&-1&-9\\1&2&-1&2}
  \]
  \[
    \xrightarrow{F_3\leftarrow F_3-F_1}
    \mwaug{ccc|c}{1&1&1&6\\0&-3&-1&-9\\0&1&-2&-4}
    \xrightarrow{F_2\leftrightarrow F_3}
    \mwaug{ccc|c}{1&1&1&6\\0&1&-2&-4\\0&-3&-1&-9}
  \]
  \[
    \xrightarrow{F_3\leftarrow F_3+3F_2}
    \mwaug{ccc|c}{1&1&1&6\\0&1&-2&-4\\0&0&-7&-21}
    \xrightarrow{F_3\leftarrow-\frac17F_3}
    \mwaug{ccc|c}{1&1&1&6\\0&1&-2&-4\\0&0&1&3}.
  \]
  Sustitución hacia atrás:
  \[
    z=3,\qquad y-2(3)=-4\Rightarrow y=2,\qquad x+2+3=6\Rightarrow x=1.
  \]
  Solución: $(x,y,z)=(1,2,3)$.
\end{cuadrosolucion}

\begin{cuadrosolucion}
  \small
  \textbf{Ejemplo 7 (sistema $3\times3$ por Gauss-Jordan).} Resolver
  \[
    \begin{cases}
      x+2y+z=5,\\
      2x+y-z=1,\\
      x-y+2z=4.
    \end{cases}
  \]

  \medskip
  \textbf{Solución.} Primero se elimina por debajo de cada pivote:
  \[
    \mwaug{ccc|c}{1&2&1&5\\2&1&-1&1\\1&-1&2&4}
    \xrightarrow{F_2\leftarrow F_2-2F_1}
    \mwaug{ccc|c}{1&2&1&5\\0&-3&-3&-9\\1&-1&2&4}
  \]
  \[
    \xrightarrow{F_3\leftarrow F_3-F_1}
    \mwaug{ccc|c}{1&2&1&5\\0&-3&-3&-9\\0&-3&1&-1}
    \xrightarrow{F_2\leftarrow-\frac13F_2}
    \mwaug{ccc|c}{1&2&1&5\\0&1&1&3\\0&-3&1&-1}
  \]
  \[
    \xrightarrow{F_3\leftarrow F_3+3F_2}
    \mwaug{ccc|c}{1&2&1&5\\0&1&1&3\\0&0&4&8}
    \xrightarrow{F_3\leftarrow\frac14F_3}
    \mwaug{ccc|c}{1&2&1&5\\0&1&1&3\\0&0&1&2}.
  \]
  Ahora se elimina también \emph{por encima} de los pivotes:
  \[
    \xrightarrow{F_2\leftarrow F_2-F_3}
    \mwaug{ccc|c}{1&2&1&5\\0&1&0&1\\0&0&1&2}
    \xrightarrow{F_1\leftarrow F_1-F_3}
    \mwaug{ccc|c}{1&2&0&3\\0&1&0&1\\0&0&1&2}
  \]
  \[
    \xrightarrow{F_1\leftarrow F_1-2F_2}
    \mwaug{ccc|c}{1&0&0&1\\0&1&0&1\\0&0&1&2}.
  \]
  La izquierda es $I_3$, así que $(x,y,z)=(1,1,2)$.
\end{cuadrosolucion}

\begin{center}
\small
\begin{tabular}{@{}p{0.30\linewidth} p{0.30\linewidth} p{0.32\linewidth}@{}}
  \toprule
  \textbf{Aspecto} & \textbf{Gauss} & \textbf{Gauss-Jordan} \\
  \midrule
  Forma final & Escalonada & Escalonada reducida \\
  Trabajo extra & Sustitución hacia atrás & Más operaciones de fila \\
  Lectura & Requiere despejar & Directa \\
  Errores típicos & Signos en la sustitución & Muchas operaciones acumuladas \\
  Cuándo conviene & Sistemas grandes & Sistemas pequeños o matrices inversas \\
  \bottomrule
\end{tabular}
\end{center}

\begin{cuadroteoria}
  \small
  \textbf{Cómo interpretar la última fila.} Al terminar la eliminación, la
  última fila de la matriz aumentada puede verse así:
  \begin{center}
  \begin{tabular}{@{}lll@{}}
    \toprule
    \textbf{Última fila} & \textbf{Significado} & \textbf{Solución} \\
    \midrule
    $[\,0\ \cdots\ 0\mid c\,]$, $c\neq0$ & $0=c$, falso & No existe \\
    $[\,0\ \cdots\ 0\mid 0\,]$ & Ecuación redundante & Infinitas \\
    $[\,0\ \cdots\ 1\mid c\,]$ & Última incógnita despejada & Única \\
    \bottomrule
  \end{tabular}
  \end{center}
\end{cuadroteoria}

\begin{cuadroadvertencia}
  \small
  \textbf{Nunca dividas por cero.} Si el pivote es $0$, \textbf{intercambia
  filas} antes de dividir. Además, no olvides que toda operación afecta a
  \textbf{toda la fila}, incluida la columna de las constantes. Cuando
  Gauss-Jordan introduce fracciones, conviene mantenerlas simplificadas en
  lugar de pasar a decimales: los errores de redondeo pueden hacer que una
  solución exacta parezca incorrecta.
\end{cuadroadvertencia}

\begin{cuadroextra}
  \small
  \textbf{Generalización.} Los mismos pasos funcionan para sistemas
  $3\times3$ o mayores: se busca un pivote en cada columna, se anula por
  debajo (Gauss) y, si se desea, también por encima (Gauss-Jordan). Un
  sistema $n\times n$ con solución única termina con la identidad $I_n$ a la
  izquierda y el vector solución a la derecha. En esta sección trabajamos
  ejemplos $2\times2$ (para ver el mecanismo con claridad) y $3\times3$ (el
  tamaño que predomina en el banco de ejercicios); los casos $2\times3$ y
  $3\times2$ se interpretan con la misma tabla de lectura.
\end{cuadroextra}

% ============================================================
\section{Determinantes}
% ============================================================

El determinante es un número asociado a una matriz \textbf{cuadrada}. Resume
mucha información: si vale $0$, la matriz es \emph{singular} y el sistema
asociado no tiene solución única; si es distinto de $0$, hay solución única.

% ------------------------------------------------------------
\subsection{Determinante de una matriz $2\times2$}
% ------------------------------------------------------------

\begin{cuadroteoria}
  \small
  \textbf{Definición.} Para
  $A=\mwmat{a&b\\c&d}$,
  \[
    \det(A)=|A|=\mwdet{a&b\\c&d}=ad-bc.
  \]
  Es el producto de la diagonal principal menos el producto de la diagonal
  secundaria.
\end{cuadroteoria}

\begin{cuadrosolucion}
  \small
  \textbf{Ejemplo 1.}
  \[
    \mwdet{3&4\\2&5}=3\cdot5-4\cdot2=15-8=7.
  \]
  \textbf{Ejemplo 2.}
  \[
    \mwdet{-2&5\\3&-4}=(-2)(-4)-5\cdot3=8-15=-7.
  \]
  \textbf{Ejemplo 3.}
  \[
    \mwdet{\tfrac12&2\\[2pt]3&\tfrac13}
    =\tfrac12\cdot\tfrac13-2\cdot3
    =\tfrac16-6=-\tfrac{35}{6}.
  \]
\end{cuadrosolucion}

\begin{cuadrosolucion}
  \small
  \textbf{Ejemplo 4 (determinante con incógnita).} Hallar $x$ si
  \[
    \mwdet{x&2\\3&x}=10.
  \]
  \textbf{Solución.}
  \[
    x\cdot x-2\cdot3=10
    \;\Rightarrow\;
    x^2-6=10
    \;\Rightarrow\;
    x^2=16
    \;\Rightarrow\;
    x=\pm4.
  \]
\end{cuadrosolucion}

% ------------------------------------------------------------
\subsection{Determinante de una matriz $3\times3$: regla de Sarrus}
% ------------------------------------------------------------

La regla de Sarrus calcula el determinante de una matriz $3\times3$ sumando
tres productos «hacia abajo» y restando tres productos «hacia arriba». Se
puede organizar de dos formas equivalentes.

\begin{cuadroteoria}
  \small
  \textbf{Sarrus horizontal.} Se repiten las dos primeras \textbf{columnas}
  a la derecha de la matriz:
  \[
    \det(A)=
    \underbrace{
    \begin{array}{ccc|cc}
      a_{11}&a_{12}&a_{13}&a_{11}&a_{12}\\
      a_{21}&a_{22}&a_{23}&a_{21}&a_{22}\\
      a_{31}&a_{32}&a_{33}&a_{31}&a_{32}
    \end{array}}_{\text{matriz ampliada}}
  \]
  \[
    \det(A)=(a_{11}a_{22}a_{33}+a_{12}a_{23}a_{31}+a_{13}a_{21}a_{32})
    -(a_{31}a_{22}a_{13}+a_{32}a_{23}a_{11}+a_{33}a_{21}a_{12}).
  \]
  \textbf{Sarrus vertical.} Se repiten las dos primeras \textbf{filas}
  debajo de la matriz y se trazan las mismas seis diagonales.
\end{cuadroteoria}

\begin{center}
\begin{tikzpicture}[baseline=(S.center)]
  \matrix (S) [matrix of math nodes,
               nodes={draw=mwprimario!40, fill=none, minimum width=1.02cm,
                      minimum height=0.72cm, inner sep=1pt, font=\scriptsize},
               row sep=-\pgflinewidth, column sep=-\pgflinewidth,
               col 4/.style={nodes={draw=mwgris!35, text=mwgris}},
               col 5/.style={nodes={draw=mwgris!35, text=mwgris}}] {
    2 & -1 & 3 & 2 & -1\\
    0 & 4 & 1 & 0 & 4\\
    5 & 2 & -2 & 5 & 2\\
  };
  \begin{scope}[on background layer]
    \draw[mwprimario, line width=4pt, opacity=0.22, line cap=round]
      (S-1-1.center) -- (S-2-2.center) -- (S-3-3.center);
    \draw[mwprimario, line width=4pt, opacity=0.22, line cap=round]
      (S-1-2.center) -- (S-2-3.center) -- (S-3-4.center);
    \draw[mwprimario, line width=4pt, opacity=0.22, line cap=round]
      (S-1-3.center) -- (S-2-4.center) -- (S-3-5.center);
    \draw[mwmagenta, line width=4pt, opacity=0.20, line cap=round]
      (S-3-1.center) -- (S-2-2.center) -- (S-1-3.center);
    \draw[mwmagenta, line width=4pt, opacity=0.20, line cap=round]
      (S-3-2.center) -- (S-2-3.center) -- (S-1-4.center);
    \draw[mwmagenta, line width=4pt, opacity=0.20, line cap=round]
      (S-3-3.center) -- (S-2-4.center) -- (S-1-5.center);
  \end{scope}
  \node[below=0.55cm of S-3-3, font=\scriptsize, mwprimario]
    {$\searrow$ productos positivos};
  \node[below=0.55cm of S-3-3, font=\scriptsize, mwmagenta, yshift=-0.42cm]
    {$\nearrow$ productos negativos};
\end{tikzpicture}
\end{center}

\begin{cuadroteoria}
  \small
  \textbf{Aviso importante.} La regla de Sarrus \textbf{solo} vale para
  matrices $3\times3$. No existe una «regla de Sarrus» para $4\times4$: para
  esas matrices hay que usar cofactores (o reducción a forma triangular).
\end{cuadroteoria}

\begin{cuadrosolucion}
  \small
  \textbf{Ejemplo 5 (Sarrus, con la matriz de la figura).} Calcular
  \[
    \det(A)=\mwdet{2&-1&3\\0&4&1\\5&2&-2}.
  \]

  \medskip
  \textbf{Solución con Sarrus horizontal.} Se repiten las dos primeras
  columnas:
  \[
    \begin{array}{ccc|cc}
      2&-1&3&2&-1\\
      0&4&1&0&4\\
      5&2&-2&5&2
    \end{array}
  \]
  Productos positivos ($\searrow$):
  \[
    (2)(4)(-2)+(-1)(1)(5)+(3)(0)(2)
    =-16-5+0=-21.
  \]
  Productos negativos ($\nearrow$):
  \[
    (5)(4)(3)+(2)(1)(2)+(-2)(0)(-1)
    =60+4+0=64.
  \]
  \[
    \det(A)=-21-64=-85.
  \]
\end{cuadrosolucion}

\begin{cuadrosolucion}
  \small
  \textbf{Ejemplo 6 (Sarrus vertical, comprobación).} Con la misma matriz,
  se repiten las dos primeras filas debajo:
  \[
    \begin{array}{ccc}
      2&-1&3\\
      0&4&1\\
      5&2&-2\\
      \hline
      2&-1&3\\
      0&4&1
    \end{array}
  \]
  Diagonales descendentes:
  $(2)(4)(-2)+(-1)(1)(5)+(3)(0)(2)=-21$.
  \quad Diagonales ascendentes:
  $(5)(4)(3)+(2)(1)(2)+(-2)(0)(-1)=64$.
  \[
    \det(A)=-21-64=-85.
  \]
  Ambos procedimientos coinciden, como debe ser.
\end{cuadrosolucion}

\begin{cuadrosolucion}
  \small
  \textbf{Ejemplo 7 (determinante nulo).} Calcular
  \[
    \mwdet{1&2&3\\4&5&6\\7&8&9}.
  \]
  \textbf{Solución.} Por Sarrus:
  \[
    (1)(5)(9)+(2)(6)(7)+(3)(4)(8)
    -(7)(5)(3)-(8)(6)(1)-(9)(4)(2)
  \]
  \[
    =45+84+96-(105+48+72)=225-225=0.
  \]
  Un determinante $0$ indica que las filas son \textbf{linealmente
  dependientes}: la tercera fila es el doble de la segunda menos la primera.
\end{cuadrosolucion}

% ------------------------------------------------------------
\subsection{Determinante de $4\times4$ o mayor: cofactores}
% ------------------------------------------------------------

\begin{cuadroteoria}
  \small
  \textbf{Menor y cofactor.} Para una matriz cuadrada $A$:
  \begin{itemize}[leftmargin=1.2em, nosep]
    \item el \textbf{menor} $M_{ij}$ es el determinante de la matriz que
          queda al \textbf{eliminar la fila $i$ y la columna $j$};
    \item el \textbf{cofactor} es
          \[
            C_{ij}=(-1)^{i+j}M_{ij}.
          \]
  \end{itemize}
  \textbf{Desarrollo por cofactores.} Para cualquier fila $i$,
  \[
    \det(A)=\sum_{j=1}^{n}a_{ij}C_{ij}
    =a_{i1}C_{i1}+a_{i2}C_{i2}+\cdots+a_{in}C_{in},
  \]
  y lo mismo para cualquier columna $j$:
  \[
    \det(A)=\sum_{i=1}^{n}a_{ij}C_{ij}.
  \]
\end{cuadroteoria}

\begin{center}
\small
\begin{tabular}{@{}cccc@{}}
  \toprule
  $+$ & $-$ & $+$ & $-$ \\
  $-$ & $+$ & $-$ & $+$ \\
  $+$ & $-$ & $+$ & $-$ \\
  $-$ & $+$ & $-$ & $+$ \\
  \bottomrule
\end{tabular}
\qquad
\text{tablero de signos } $(-1)^{i+j}$
\end{center}

\begin{center}
\begin{tikzpicture}[baseline=(M.center)]
  \matrix (M) [matrix of math nodes,
               nodes={draw=mwprimario!40, fill=white, minimum width=1.0cm,
                      minimum height=0.72cm, inner sep=1pt, font=\scriptsize},
               row sep=-\pgflinewidth, column sep=-\pgflinewidth,
               row 1/.style={nodes={fill=mwgris!16, text=mwgris!60,
                                    draw=mwgris!30}},
               column 1/.style={nodes={fill=mwgris!16, text=mwgris!60,
                                       draw=mwgris!30}}] {
    2 & 0 & 0 & 0\\
    3 & 1 & 2 & 0\\
    4 & 0 & 3 & 1\\
    5 & 2 & 1 & 2\\
  };
  \draw[mwvioleta, line width=1.2pt, rounded corners=1pt]
    (M-1-1.north west) rectangle (M-1-1.south east);
  \matrix (N) [matrix of math nodes,
               nodes={draw=mwprimario!40, fill=white, minimum width=1.0cm,
                      minimum height=0.72cm, inner sep=1pt, font=\scriptsize},
               right=1.6cm of M,
               row sep=-\pgflinewidth, column sep=-\pgflinewidth] {
    1 & 2 & 0\\
    0 & 3 & 1\\
    2 & 1 & 2\\
  };
  \node[right=0.3cm of M, font=\large] {$\longrightarrow$};
  \node[above=0.3cm of N, font=\scriptsize, mwmagenta] {$M_{11}$ (menor)};
  \node[below=0.35cm of M, font=\scriptsize, mwgris, align=center]
    {se eliminan la fila $1$\\ y la columna $1$};
  \node[below=0.35cm of N, font=\scriptsize, mwgris]
    {matriz $3\times3$ restante};
\end{tikzpicture}
\end{center}

\begin{cuadrosolucion}
  \small
  \textbf{Ejemplo 8 ($4\times4$ por cofactores).} Calcular
  \[
    \det(M)=\mwdet{2&0&0&0\\3&1&2&0\\4&0&3&1\\5&2&1&2}.
  \]

  \medskip
  \textbf{Solución.} La primera fila tiene tres ceros: conviene desarrollar
  por ella. Como $a_{11}=2$ y los demás $a_{1j}=0$,
  \[
    \det(M)=2\cdot(-1)^{1+1}M_{11}
    =2\mwdet{1&2&0\\0&3&1\\2&1&2}.
  \]
  El menor $3\times3$ se calcula por Sarrus:
  \[
    \mwdet{1&2&0\\0&3&1\\2&1&2}
    =(1)(3)(2)+(2)(1)(2)+(0)(0)(1)
    -(2)(3)(0)-(1)(1)(1)-(2)(0)(2)
  \]
  \[
    =6+4+0-(0+1+0)=9.
  \]
  Por lo tanto,
  \[
    \det(M)=2\cdot9=18.
  \]
\end{cuadrosolucion}

\begin{cuadropasos}
  \small
  \textbf{Cómo elegir la fila o columna para desarrollar:}
  \begin{enumerate}[leftmargin=1.4em, nosep]
    \item busca la fila o columna con \textbf{más ceros};
    \item si hay empate, elige la que tenga números más pequeños;
    \item recuerda que cada término lleva el signo $(-1)^{i+j}$;
    \item los menores se calculan con Sarrus si son $3\times3$, o se vuelve
          a aplicar cofactores si son mayores.
  \end{enumerate}
  Para $5\times5$ o más, el procedimiento es \textbf{recursivo}: cada
  cofactor se reduce a un determinante de orden una unidad menor, hasta
  llegar a $3\times3$ o $2\times2$.
\end{cuadropasos}

\begin{cuadroteoria}
  \small
  \textbf{Propiedades de los determinantes.} Para matrices cuadradas del
  mismo orden:
  \begin{enumerate}[leftmargin=1.4em, nosep]
    \item $\det(I)=1$ y $\det(O)=0$.
    \item Si $A$ es triangular (superior o inferior), $\det(A)$ es el
          \textbf{producto de la diagonal}.
    \item Al \textbf{intercambiar dos filas}, el determinante cambia de
          signo.
    \item Al \textbf{multiplicar una fila por $k$}, el determinante queda
          multiplicado por $k$. En particular,
          $\det(kA)=k^n\det(A)$.
    \item Al \textbf{sumar a una fila un múltiplo de otra}, el determinante
          \textbf{no cambia}.
    \item Si una fila o columna es nula, o es múltiplo de otra,
          $\det(A)=0$.
    \item $\det(AB)=\det(A)\det(B)$ y $\det(A^{\top})=\det(A)$.
  \end{enumerate}
\end{cuadroteoria}

\begin{cuadroadvertencia}
  \small
  \textbf{Los dos errores más caros:}
  \begin{itemize}[leftmargin=1.2em, nosep]
    \item \textbf{olvidar el signo del cofactor}: el tablero
          $+,-,+,-$ alterna por filas y columnas;
    \item \textbf{eliminar la fila equivocada}: para $M_{ij}$ se quita la
          fila $i$ y la columna $j$ del elemento elegido, no las primeras.
  \end{itemize}
  Además, el determinante \textbf{solo existe para matrices cuadradas}: no
  se calcula el determinante de una matriz $2\times3$.
\end{cuadroadvertencia}

% ============================================================
\section{Resolución de sistemas por la regla de Cramer}
% ============================================================

\begin{cuadroteoria}
  \small
  \textbf{Condición.} La regla de Cramer se aplica a un sistema
  \textbf{cuadrado} (igual número de ecuaciones que de incógnitas) cuya
  matriz de coeficientes tenga \textbf{determinante distinto de cero}:
  \[
    \det(A)\neq0.
  \]
  En ese caso el sistema es compatible determinado y su solución es
  \[
    x_j=\frac{\det(A_j)}{\det(A)},\qquad j=1,2,\dots,n,
  \]
  donde $A_j$ es la matriz que resulta de sustituir la \textbf{columna $j$}
  de $A$ por el vector de constantes $\mathbf{b}$.
\end{cuadroteoria}

Para un sistema $2\times2$,
\[
  \begin{cases}
    a_1x+b_1y=c_1,\\
    a_2x+b_2y=c_2,
  \end{cases}
\]
las fórmulas son explícitas:
\[
  x=\frac{\mwdet{c_1&b_1\\c_2&b_2}}{\mwdet{a_1&b_1\\a_2&b_2}},
  \qquad
  y=\frac{\mwdet{a_1&c_1\\a_2&c_2}}{\mwdet{a_1&b_1\\a_2&b_2}}.
\]

\begin{cuadrosolucion}
  \small
  \textbf{Ejemplo 1.} Resolver por Cramer
  \[
    \begin{cases}
      2x+3y=8,\\
      x+4y=9.
    \end{cases}
  \]

  \medskip
  \textbf{Solución.} Primero el determinante del sistema:
  \[
    \det(A)=\mwdet{2&3\\1&4}=2\cdot4-3\cdot1=8-3=5\neq0.
  \]
  Como no es cero, Cramer se puede aplicar:
  \[
    x=\frac{\mwdet{8&3\\9&4}}{\det(A)}
     =\frac{8\cdot4-3\cdot9}{5}
     =\frac{32-27}{5}=\frac{5}{5}=1,
  \]
  \[
    y=\frac{\mwdet{2&8\\1&9}}{\det(A)}
     =\frac{2\cdot9-8\cdot1}{5}
     =\frac{18-8}{5}=\frac{10}{5}=2.
  \]
  Solución: $(x,y)=(1,2)$.
\end{cuadrosolucion}

\begin{cuadrosolucion}
  \small
  \textbf{Ejemplo 2.} Resolver por Cramer
  \[
    \begin{cases}
      x+2y=7,\\
      3x-y=7.
    \end{cases}
  \]

  \medskip
  \textbf{Solución.}
  \[
    \det(A)=\mwdet{1&2\\3&-1}=(1)(-1)-(2)(3)=-1-6=-7.
  \]
  \[
    x=\frac{\mwdet{7&2\\7&-1}}{-7}
     =\frac{-7-14}{-7}=\frac{-21}{-7}=3,
    \qquad
    y=\frac{\mwdet{1&7\\3&7}}{-7}
     =\frac{7-21}{-7}=\frac{-14}{-7}=2.
  \]
  Solución: $(x,y)=(3,2)$.
\end{cuadrosolucion}

\begin{cuadrosolucion}
  \small
  \textbf{Ejemplo 3 (determinante nulo).} Intentar resolver por Cramer
  \[
    \begin{cases}
      2x+4y=6,\\
      x+2y=5.
    \end{cases}
  \]

  \medskip
  \textbf{Solución.}
  \[
    \det(A)=\mwdet{2&4\\1&2}=4-4=0.
  \]
  Cramer \textbf{no se puede aplicar}: la fórmula exige dividir entre
  $\det(A)$, que es cero. Además el sistema es \textbf{incompatible}: la
  primera ecuación es $2(x+2y)=6$, es decir $x+2y=3$, mientras que la
  segunda afirma $x+2y=5$. No existe solución.
\end{cuadrosolucion}

\begin{cuadrosolucion}
  \small
  \textbf{Ejemplo 4 (sistema $3\times3$ por Cramer).} Resolver
  \[
    \begin{cases}
      x+y+z=6,\\
      2x+3y+z=13,\\
      x-y+2z=3.
    \end{cases}
  \]

  \medskip
  \textbf{Solución.} Determinante del sistema, por Sarrus:
  \[
    \det(A)=\mwdet{1&1&1\\2&3&1\\1&-1&2}
    =(6+1-2)-(3-1+4)=5-6=-1\neq0.
  \]
  Ahora se sustituye cada columna por $\mathbf{b}=(6,13,3)$:
  \[
    \det(A_x)=\mwdet{6&1&1\\13&3&1\\3&-1&2}=-3,
    \quad
    \det(A_y)=\mwdet{1&6&1\\2&13&1\\1&3&2}=-2,
  \]
  \[
    \det(A_z)=\mwdet{1&1&6\\2&3&13\\1&-1&3}=-1.
  \]
  Por lo tanto,
  \[
    x=\frac{-3}{-1}=3,\qquad
    y=\frac{-2}{-1}=2,\qquad
    z=\frac{-1}{-1}=1.
  \]
  Solución: $(x,y,z)=(3,2,1)$. Comprobación: $3+2+1=6$,
  $2(3)+3(2)+1=13$ y $3-2+2(1)=3$. \checkmark
\end{cuadrosolucion}

\begin{cuadroteoria}
  \small
  \textbf{Qué significa $\det(A)=0$.} Si el determinante del sistema es
  cero, la regla de Cramer no decide. Hay que analizar el sistema con Gauss:
  \begin{itemize}[leftmargin=1.2em, nosep]
    \item si al reducir aparece una fila $[\,0\ \cdots\ 0\mid c\,]$ con
          $c\neq0$, el sistema es \textbf{incompatible};
    \item si aparece una fila $[\,0\ \cdots\ 0\mid 0\,]$, el sistema tiene
          \textbf{infinitas soluciones} (al menos una incógnita libre).
  \end{itemize}
\end{cuadroteoria}

\begin{center}
\small
\begin{tabular}{@{}p{0.25\linewidth} p{0.34\linewidth} p{0.33\linewidth}@{}}
  \toprule
  \textbf{Aspecto} & \textbf{Cramer} & \textbf{Gauss / Gauss-Jordan} \\
  \midrule
  Requisito & Sistema cuadrado con $\det(A)\neq0$ & Cualquier sistema \\
  Trabajo & Un determinante por incógnita & Una sola reducción \\
  Ventaja & Fórmula directa y elegante & Funciona siempre \\
  Desventaja & No sirve si $\det(A)=0$ & Más pasos y más cuidado con signos \\
  Cuándo conviene & Sistemas $2\times2$ o $3\times3$ & Sistemas grandes \\
  \bottomrule
\end{tabular}
\end{center}

\begin{cuadroextra}
  \small
  \textbf{Las mismas fórmulas en $3\times3$.} Si
  $A=\mwmat{a_{11}&a_{12}&a_{13}\\a_{21}&a_{22}&a_{23}\\a_{31}&a_{32}&a_{33}}$,
  \[
    x=\frac{\det(A_x)}{\det(A)},\quad
    y=\frac{\det(A_y)}{\det(A)},\quad
    z=\frac{\det(A_z)}{\det(A)},
  \]
  donde $A_x$, $A_y$ y $A_z$ se obtienen sustituyendo la \emph{primera},
  \emph{segunda} y \emph{tercera} columna por el vector de constantes
  $\mathbf{b}$, y cada determinante se calcula con Sarrus. Cramer sigue
  exigiendo $\det(A)\neq0$; con tres incógnitas hay que calcular
  \textbf{cuatro} determinantes de $3\times3$, razón por la cual Gauss suele
  ser más económico.
\end{cuadroextra}

\begin{cuadroadvertencia}
  \small
  \textbf{Dos despistes típicos en Cramer:}
  \begin{itemize}[leftmargin=1.2em, nosep]
    \item sustituir la \textbf{columna equivocada}: para $x$ se reemplaza la
          \emph{primera} columna; para $y$, la \emph{segunda};
    \item invertir numerador y denominador, o cambiar el orden de las
          columnas al calcular un determinante (eso cambia el signo).
  \end{itemize}
  Antes de dividir, comprueba que $\det(A)\neq0$. Y al terminar,
  \textbf{verifica} la solución en el sistema original.
\end{cuadroadvertencia}


% ============================================================
\section{Matriz inversa}
% ============================================================

Hasta aquí hemos \emph{sumado}, \emph{multiplicado} y \emph{resuelto} sistemas.
Falta la operación que, en matrices, hace el papel de la división: la
\textbf{inversa}. No todas las matrices la tienen, pero cuando existe permite
resolver $A\mathbf{x}=\mathbf{b}$ con un solo producto:
$\mathbf{x}=A^{-1}\mathbf{b}$.

\begin{cuadroteoria}
  \small
  \textbf{Definición.} Sea $A$ una matriz \textbf{cuadrada} de orden $n$. Se
  dice que $A$ es \textbf{invertible} (o \emph{no singular}) si existe una
  matriz $B$ del mismo orden tal que
  \[
    AB=BA=I_n .
  \]
  Esa matriz $B$, si existe, es \textbf{única} y se llama \textbf{inversa} de
  $A$: se escribe $B=A^{-1}$. Así,
  \[
    AA^{-1}=A^{-1}A=I_n .
  \]
  Si no existe tal matriz, $A$ es \textbf{singular} y no tiene inversa.
\end{cuadroteoria}

\begin{cuadroteoria}
  \small
  \textbf{Cuándo existe: el criterio de equivalencia.} Para una matriz
  cuadrada $A$ son equivalentes:
  \begin{enumerate}[leftmargin=1.4em, nosep]
    \item $A$ es invertible;
    \item $\det(A)\neq0$;
    \item la forma escalonada reducida de $A$ es $I_n$;
    \item la matriz aumentada $[A\mid I]$ es equivalente por filas a
          $[I\mid A^{-1}]$.
  \end{enumerate}
  La última es la que usaremos para \emph{calcular} la inversa: si al reducir
  la mitad izquierda aparece $I$, la mitad derecha \emph{es} $A^{-1}$; si
  aparece una fila de ceros a la izquierda, la inversa \textbf{no existe}.
\end{cuadroteoria}

\begin{cuadropasos}
  \small
  \textbf{Método de Gauss-Jordan para hallar $A^{-1}$:}
  \begin{enumerate}[leftmargin=1.4em, nosep]
    \item Escribir la matriz aumentada $[A\mid I]$.
    \item Aplicar operaciones elementales de fila \textbf{a toda la fila},
          incluida la parte de la identidad: cada operación que se hace a la
          izquierda se hace también a la derecha.
    \item Continuar hasta que la izquierda sea $I_n$.
    \item Leer la inversa en la mitad derecha: $[I\mid A^{-1}]$.
    \item Si en la mitad izquierda aparece una fila nula, detenerse:
          $A$ no es invertible.
  \end{enumerate}
\end{cuadropasos}

\begin{cuadrosolucion}
  \small
  \textbf{Ejemplo 1 ($2\times2$).} Hallar la inversa de
  $A=\mwmat{3&1\\2&1}$.

  \medskip
  \textbf{Solución.} Se parte de $[A\mid I]$ y se reduce:
  \[
    \mwaug{cc|cc}{3&1&1&0\\2&1&0&1}
    \xrightarrow{F_1\leftarrow F_1-F_2}
    \mwaug{cc|cc}{1&0&1&-1\\2&1&0&1}
  \]
  \[
    \xrightarrow{F_2\leftarrow F_2-2F_1}
    \mwaug{cc|cc}{1&0&1&-1\\0&1&-2&3}.
  \]
  La izquierda ya es $I_2$, así que
  \[
    A^{-1}=\mwmat{1&-1\\-2&3}.
  \]
  \textbf{Comprobación} (siempre conviene hacerla):
  \[
    AA^{-1}=\mwmat{3&1\\2&1}\mwmat{1&-1\\-2&3}
    =\mwmat{1&0\\0&1}=I_2. \checkmark
  \]
\end{cuadrosolucion}

\begin{cuadrosolucion}
  \small
  \textbf{Ejemplo 2 ($2\times2$ con determinante negativo).} Hallar la inversa
  de $A=\mwmat{4&3\\3&2}$.

  \medskip
  \textbf{Solución.}
  \[
    \mwaug{cc|cc}{4&3&1&0\\3&2&0&1}
    \xrightarrow{F_1\leftarrow F_1-F_2}
    \mwaug{cc|cc}{1&1&1&-1\\3&2&0&1}
    \xrightarrow{F_2\leftarrow F_2-3F_1}
    \mwaug{cc|cc}{1&1&1&-1\\0&-1&-3&4}
  \]
  \[
    \xrightarrow{F_2\leftarrow -F_2}
    \mwaug{cc|cc}{1&1&1&-1\\0&1&3&-4}
    \xrightarrow{F_1\leftarrow F_1-F_2}
    \mwaug{cc|cc}{1&0&-2&3\\0&1&3&-4}.
  \]
  Por lo tanto,
  \[
    A^{-1}=\mwmat{-2&3\\3&-4}.
  \]
  \textbf{Comprobación:} $4(-2)+3(3)=1$ y $3(3)+2(-4)=1$ en la diagonal,
  mientras que los productos cruzados dan $0$. \checkmark
\end{cuadrosolucion}

\begin{cuadroteoria}
  \small
  \textbf{Atajo para $2\times2$.} Si $A=\mwmat{a&b\\c&d}$ y
  $\det(A)=ad-bc\neq0$, entonces
  \[
    A^{-1}=\frac{1}{ad-bc}\mwmat{d&-b\\-c&a}.
  \]
  Se \textbf{intercambian} $a$ y $d$, se \textbf{cambian de signo} $b$ y $c$, y
  se divide todo entre el determinante. Es un atajo cómodo para $2\times2$;
  para $3\times3$ o mayor no hay fórmula corta y se usa
  Gauss-Jordan sobre $[A\mid I]$.
\end{cuadroteoria}

\begin{cuadrosolucion}
  \small
  \textbf{Ejemplo 3 ($3\times3$).} Hallar la inversa de
  $A=\mwmat{1&2&3\\0&1&4\\5&6&0}$.

  \medskip
  \textbf{Solución.} Se aumenta con $I_3$ y se elimina columna por columna:
  \[
    \mwaug{ccc|ccc}{1&2&3&1&0&0\\
                    0&1&4&0&1&0\\
                    5&6&0&0&0&1}
    \xrightarrow{F_3\leftarrow F_3-5F_1}
    \mwaug{ccc|ccc}{1&2&3&1&0&0\\
                    0&1&4&0&1&0\\
                    0&-4&-15&-5&0&1}
  \]
  \[
    \xrightarrow{F_3\leftarrow F_3+4F_2}
    \mwaug{ccc|ccc}{1&2&3&1&0&0\\
                    0&1&4&0&1&0\\
                    0&0&1&-5&4&1}.
  \]
  Ahora se anula hacia arriba. Primero la tercera columna:
  \[
    \xrightarrow{F_2\leftarrow F_2-4F_3}
    \mwaug{ccc|ccc}{1&2&3&1&0&0\\
                    0&1&0&20&-15&-4\\
                    0&0&1&-5&4&1}
  \]
  \[
    \xrightarrow{F_1\leftarrow F_1-3F_3}
    \mwaug{ccc|ccc}{1&2&0&16&-12&-3\\
                    0&1&0&20&-15&-4\\
                    0&0&1&-5&4&1}
    \xrightarrow{F_1\leftarrow F_1-2F_2}
    \mwaug{ccc|ccc}{1&0&0&-24&18&5\\
                    0&1&0&20&-15&-4\\
                    0&0&1&-5&4&1}.
  \]
  Por lo tanto,
  \[
    A^{-1}=\mwmat{-24&18&5\\20&-15&-4\\-5&4&1}.
  \]
  \textbf{Comprobación de una entrada:} la primera fila de $A$ por la primera
  columna de $A^{-1}$ da $1(-24)+2(20)+3(-5)=-24+40-15=1$, como debe ser.
  \checkmark
\end{cuadrosolucion}

\begin{cuadrosolucion}
  \small
  \textbf{Ejemplo 4 (matriz sin inversa).} Intentar hallar la inversa de
  $A=\mwmat{1&2\\2&4}$.

  \medskip
  \textbf{Solución.}
  \[
    \mwaug{cc|cc}{1&2&1&0\\2&4&0&1}
    \xrightarrow{F_2\leftarrow F_2-2F_1}
    \mwaug{cc|cc}{1&2&1&0\\0&0&-2&1}.
  \]
  En la mitad izquierda apareció una \textbf{fila nula}: es imposible
  convertirla en $I_2$. Por lo tanto $A$ \textbf{no es invertible}. Coincide
  con el criterio del determinante: $\det(A)=1\cdot4-2\cdot2=0$.
\end{cuadrosolucion}

\begin{cuadrosolucion}
  \small
  \textbf{Ejemplo 5 (resolver un sistema con la inversa).} Resolver
  \[
    \begin{cases}
      3x+y=5,\\
      2x+y=3
    \end{cases}
    \qquad\text{usando }A^{-1}.
  \]

  \medskip
  \textbf{Solución.} Del ejemplo 1, $A^{-1}=\mwmat{1&-1\\-2&3}$. Como
  $A\mathbf{x}=\mathbf{b}$ implica $\mathbf{x}=A^{-1}\mathbf{b}$:
  \[
    \mwmat{x\\y}=\mwmat{1&-1\\-2&3}\mwmat{5\\3}
    =\mwmat{1\cdot5-1\cdot3\\-2\cdot5+3\cdot3}
    =\mwmat{2\\-1}.
  \]
  Solución: $(x,y)=(2,-1)$. Comprobación: $3(2)+(-1)=5$ y $2(2)+(-1)=3$.
  \checkmark
\end{cuadrosolucion}

\begin{cuadroteoria}
  \small
  \textbf{Propiedades de la inversa.} Siempre que las matrices sean
  invertibles y del orden adecuado:
  \begin{enumerate}[leftmargin=1.4em, nosep]
    \item $(A^{-1})^{-1}=A$ \quad y \quad $I_n^{-1}=I_n$.
    \item $(AB)^{-1}=B^{-1}A^{-1}$ \quad (\textbf{se invierte el orden}).
    \item $(A^{\top})^{-1}=(A^{-1})^{\top}$.
    \item $(kA)^{-1}=\dfrac{1}{k}A^{-1}$, con $k\neq0$.
    \item $\det(A^{-1})=\dfrac{1}{\det(A)}$.
    \item Si $A\mathbf{x}=\mathbf{b}$ y $A$ es invertible, entonces
          $\mathbf{x}=A^{-1}\mathbf{b}$: solución única.
  \end{enumerate}
\end{cuadroteoria}

\begin{cuadroadvertencia}
  \small
  \textbf{Cinco errores que hay que evitar:}
  \begin{itemize}[leftmargin=1.2em, nosep]
    \item \textbf{Invertir entrada por entrada.} $A^{-1}$ no se obtiene
          poniendo $1/a_{ij}$; hay que reducir $[A\mid I]$.
    \item \textbf{Olvidar que solo existe para matrices cuadradas.} Una
          matriz $2\times3$ no tiene inversa (ni por izquierda ni por
          derecha en el sentido de $AA^{-1}=I$).
    \item \textbf{No comprobar el determinante.} Si $\det(A)=0$, no hay
          inversa y todo el trabajo de reducción es inútil.
    \item \textbf{Equivocar el orden en $(AB)^{-1}$:} es $B^{-1}A^{-1}$, no
          $A^{-1}B^{-1}$.
    \item \textbf{Operar solo la mitad izquierda.} Cada operación de fila
          debe aplicarse también a la identidad de la derecha.
  \end{itemize}
\end{cuadroadvertencia}

\begin{cuadroextra}
  \small
  \textbf{De dónde sale la fórmula $2\times2$.} Para
  $A=\mwmat{a&b\\c&d}$ con $\det(A)\neq0$,
  \[
    A^{-1}=\frac{1}{\det(A)}\mwmat{d&-b\\-c&a}
    =\frac{1}{\det(A)}\,\operatorname{adj}(A),
  \]
  donde $\operatorname{adj}(A)$ es la \textbf{matriz adjunta}: la traspuesta
  de la matriz de cofactores. Esta misma idea funciona en $3\times3$, pero
  exige calcular nueve cofactores; por eso, en la práctica, Gauss-Jordan
  sobre $[A\mid I]$ resulta más económico.
\end{cuadroextra}

% ============================================================
\section{Formulario resumen}
% ============================================================

\begin{cuadroteoria}
  \small
  \textbf{Definiciones y tamaños}
  \begin{itemize}[leftmargin=1.2em, nosep]
    \item Matriz $m\times n$: $m$ filas, $n$ columnas, $mn$ elementos.
    \item $A=(a_{ij})$: $i$ = fila, $j$ = columna.
    \item $A=B$ si y solo si tienen el mismo orden y $a_{ij}=b_{ij}$.
    \item Cuadrada: $m=n$. Identidad $I_n$: $1$ en la diagonal, $0$ fuera.
    \item Nula $O$: todos sus elementos son $0$.
  \end{itemize}

  \medskip
  \textbf{Operaciones}
  \begin{itemize}[leftmargin=1.2em, nosep]
    \item Escalar: $(kA)_{ij}=k\,a_{ij}$.
    \item Suma y resta: mismo orden; se opera entrada por entrada.
    \item Producto: $A_{m\times n}B_{n\times p}=C_{m\times p}$ con
          $c_{ij}=\sum_{k=1}^{n}a_{ik}b_{kj}$.
    \item $AB\neq BA$ en general; $AB=O$ no implica $A=O$ ni $B=O$.
    \item $AI=IA=A$; $A^n=AA\cdots A$ ($n$ factores).
  \end{itemize}

  \medskip
  \textbf{Sistemas}
  \begin{itemize}[leftmargin=1.2em, nosep]
    \item $A\mathbf{x}=\mathbf{b}$; aumentada $[A\mid\mathbf{b}]$.
    \item Operaciones elementales: $F_i\leftrightarrow F_j$;
          $F_i\leftarrow kF_i$ ($k\neq0$);
          $F_i\leftarrow F_i+kF_j$.
    \item Gauss: escalonada + sustitución hacia atrás.
    \item Gauss-Jordan: escalonada reducida, lectura directa.
    \item Última fila $[0\ \cdots\ 0\mid c]$, $c\neq0$: incompatible.
    \item Última fila $[0\ \cdots\ 0\mid 0]$: infinitas soluciones.
  \end{itemize}
\end{cuadroteoria}

\begin{cuadroteoria}
  \small
  \textbf{Determinantes}
  \[
    \mwdet{a&b\\c&d}=ad-bc.
  \]
  \[
    \det(A)=(a_{11}a_{22}a_{33}+a_{12}a_{23}a_{31}+a_{13}a_{21}a_{32})
    -(a_{31}a_{22}a_{13}+a_{32}a_{23}a_{11}+a_{33}a_{21}a_{12}).
  \]
  \[
    C_{ij}=(-1)^{i+j}M_{ij},
    \qquad
    \det(A)=\sum_{j=1}^{n}a_{ij}C_{ij}.
  \]
  \begin{itemize}[leftmargin=1.2em, nosep]
    \item Triangular: determinante = producto de la diagonal.
    \item Intercambiar filas: cambia el signo.
    \item Sumar un múltiplo de otra fila: no cambia el determinante.
    \item $\det(AB)=\det(A)\det(B)$; $\det(kA)=k^n\det(A)$.
  \end{itemize}

  \medskip
  \textbf{Matriz inversa} (solo si $A$ es cuadrada y $\det(A)\neq0$):
  \[
    [A\mid I]\;\sim\;[I\mid A^{-1}]
  \]
  \[
    \mwmat{a&b\\c&d}^{-1}
    =\frac{1}{ad-bc}\mwmat{d&-b\\-c&a}.
  \]
  \begin{itemize}[leftmargin=1.2em, nosep]
    \item $(AB)^{-1}=B^{-1}A^{-1}$; \quad $(A^{-1})^{-1}=A$;
          \quad $\det(A^{-1})=1/\det(A)$.
    \item Si $\det(A)=0$: $A$ es singular y \textbf{no} tiene inversa.
    \item $A\mathbf{x}=\mathbf{b}\Rightarrow\mathbf{x}=A^{-1}\mathbf{b}$.
  \end{itemize}

  \medskip
  \textbf{Cramer} (solo si $\det(A)\neq0$): $x_j=\dfrac{\det(A_j)}{\det(A)}$,
  sustituyendo la columna $j$ por $\mathbf{b}$. En $2\times2$:
  \[
    x=\frac{\mwdet{c_1&b_1\\c_2&b_2}}{\det(A)},
    \qquad
    y=\frac{\mwdet{a_1&c_1\\a_2&c_2}}{\det(A)};
  \]
  en $3\times3$ se añade $z=\det(A_z)/\det(A)$ y cada determinante se hace
  por Sarrus.
\end{cuadroteoria}

\begin{center}
\small
\begin{tabular}{@{}p{0.44\linewidth} p{0.48\linewidth}@{}}
  \toprule
  \textbf{Si el problema pide\dots} & \textbf{Conviene usar\dots} \\
  \midrule
  Sumar, restar o multiplicar por un número & Operaciones entrada por
    entrada, verificando el orden. \\
  Multiplicar dos matrices & Comprobar $(m\times n)(n\times p)$ y usar
    fila por columna. \\
  Resolver un sistema $2\times2$ o $3\times3$ & Gauss-Jordan (rápido) o
    Cramer si $\det(A)\neq0$. \\
  Decidir si un sistema tiene solución única & Calcular $\det(A)$: si es
    distinto de $0$, es única. \\
  Resolver un sistema con parámetros o sin solución única & Gauss: permite
    ver incompatibilidad e infinitas soluciones. \\
  Determinante $3\times3$ & Sarrus (horizontal o vertical). \\
  Determinante $4\times4$ o mayor & Cofactores por la fila o columna con
    más ceros. \\
  Hallar la inversa de una matriz & Reducir $[A\mid I]$ a $[I\mid A^{-1}]$
    (y comprobar que $\det(A)\neq0$). \\
  \bottomrule
\end{tabular}
\end{center}

% ============================================================
\section{Banco de ejercicios y problemas}
% ============================================================

\begin{cuadroinfo}
  \small
  \textbf{Instrucciones.} Resuelve cada ejercicio en tu cuaderno y anota el
  procedimiento, no solo el resultado. Los ejercicios están agrupados por
  tema y numerados de forma continua del $1$ al $88$. La dificultad se indica
  con estrellas: \facil\ básico, \medio\ intermedio, \dificil\ avanzado.
  La clave comentada está al final de la guía.
\end{cuadroinfo}

\mwbloquehead{Bloque A — Definición, filas, columnas y orden}{8}{0}
\begin{mwbloque}
  \item \facil\ Dada $A=\mwmat{2&-1&3\\0&5&-4}$, indicar su orden, sus filas,
        sus columnas y los elementos $a_{12}$, $a_{23}$ y $a_{21}$.
  \item \facil\ Dada $B=\mwmat{1&0\\-2&3\\4&5}$, indicar su orden, el número
        de elementos, los elementos $b_{31}$ y $b_{22}$, y escribir la
        primera columna.
  \item \facil\ Construir la matriz $C=(c_{ij})$ de orden $2\times3$ definida
        por $c_{ij}=i+2j$.
  \item \facil\ Construir la matriz $D=(d_{ij})$ de orden $3\times2$ definida
        por $d_{ij}=i-j$.
  \item \facil\ Construir la matriz $E=(e_{ij})$ de orden $2\times2$ definida
        por $e_{ij}=2i+j$.
  \item \facil\ Hallar $x$ e $y$ para que
        $\mwmat{x+1&3\\4&2y}=\mwmat{5&3\\4&10}$.
  \item \medio\ Indicar en qué casos las dos matrices pueden ser iguales:
        (a) una $2\times3$ y una $3\times2$; (b) dos $2\times2$;
        (c) una $1\times3$ y una $3\times1$.
  \item \facil\ Completar: (a) una matriz $4\times5$ tiene \rule{1cm}{0.4pt}
        elementos; (b) una $7\times1$ tiene \rule{1cm}{0.4pt};
        (c) una $n\times n$ tiene \rule{1cm}{0.4pt};
        (d) una $3\times4$ tiene \rule{1cm}{0.4pt}.
\end{mwbloque}

\mwbloquehead{Bloque B — Multiplicación por un escalar}{16}{8}
\begin{mwbloque}
  \item \facil\ Si $A=\mwmat{2&-3\\1&4}$, calcular $3A$.
  \item \facil\ Si $B=\mwmat{-1&2\\0&5\\3&-2}$, calcular $-2B$.
  \item \facil\ Si $C=\mwmat{\tfrac12&-1\\[2pt]3&\tfrac23}$, calcular $6C$.
  \item \medio\ Si $A=\mwmat{2&1\\-1&3}$ y $B=\mwmat{0&-2\\4&1}$, calcular
        $2A-3B$.
  \item \medio\ Resolver para $X$:
        $2X+\mwmat{1&2\\3&4}=\mwmat{5&6\\7&8}$.
  \item \facil\ Si $A=\mwmat{1&-2\\3&0}$, calcular $-\tfrac12A$.
  \item \medio\ Verificar que $3(A+B)=3A+3B$ para
        $A=\mwmat{1&0\\2&-1}$ y $B=\mwmat{-1&3\\0&2}$.
  \item \facil\ Hallar el escalar $k$ tal que
        $k\mwmat{2&-1\\4&6}=\mwmat{-6&3\\-12&-18}$.
\end{mwbloque}

\mwbloquehead{Bloque C — Suma y resta de matrices}{26}{16}
\begin{mwbloque}
  \item \facil\ Con $A=\mwmat{1&2\\3&4}$ y $B=\mwmat{5&-1\\0&2}$, calcular
        $A+B$ y $A-B$.
  \item \medio\ Con $A=\mwmat{2&0&-1\\3&4&5}$ y
        $B=\mwmat{-1&2&0\\1&-3&2}$, calcular $A+B$ y $A-B$.
  \item \facil\ Con $A=\mwmat{4&-2\\1&0\\3&5}$ y
        $B=\mwmat{-1&2\\0&-3\\2&1}$, calcular $A+B$.
  \item \medio\ ¿Es posible calcular $A+B$ si $A$ es $2\times3$ y $B$ es
        $3\times2$? Justificar.
  \item \medio\ Resolver para $X$:
        $X-\mwmat{2&1\\0&3}=\mwmat{4&-1\\5&2}$.
  \item \medio\ Hallar $a$ y $b$ para que
        $\mwmat{a&2\\3&b}+\mwmat{1&-2\\-3&4}=\mwmat{4&0\\0&7}$.
  \item \medio\ Con $A=\mwmat{1&0\\2&-1}$, $B=\mwmat{3&1\\-1&2}$ y
        $C=\mwmat{0&4\\1&0}$, calcular $2A-B+C$.
  \item \facil\ Verificar que $A+O=A$ para $A=\mwmat{3&-1\\2&0}$ y
        $O=\mwmat{0&0\\0&0}$.
  \item \medio\ Con $A=\mwmat{1&2\\3&4}$ y $B=\mwmat{5&-1\\0&2}$, calcular
        $(A+B)-(A-B)$ y comparar con $2B$.
  \item \facil\ Hallar $x$ e $y$ si
        $\mwmat{x&1\\2&y}-\mwmat{1&1\\2&3}=O$.
\end{mwbloque}

\mwbloquehead{Bloque D — Sistemas y matriz aumentada}{34}{26}
\begin{mwbloque}
  \item \facil\ Escribir la matriz aumentada del sistema
        $x+2y-z=3$, $2x-y+z=0$, $3x+y+2z=7$.
  \item \facil\ Escribir el sistema asociado a
        $\mwaug{ccc|c}{1&1&1&6\\2&-1&1&3\\1&2&-1&2}$ y comprobar si
        $(1,2,3)$ es solución.
  \item \facil\ Escribir en forma matricial $A\mathbf{x}=\mathbf{b}$ el
        sistema $2x-y+3z=5$, $x+4y-z=2$, $3x+y+z=8$.
  \item \facil\ Para el sistema $x+2y+3z=6$, $2x-y+z=1$, $x+y-z=2$,
        escribir $A$, $\mathbf{x}$, $\mathbf{b}$ y la matriz aumentada.
  \item \dificil\ En el sistema $x+y+z=3$, $2x-y+z=0$ (caso $2\times3$),
        escribir la matriz aumentada y dar dos soluciones distintas.
  \item \medio\ Explicar por qué el sistema $x+y+z=1$, $x+y+z=2$ (caso
        $2\times3$) no tiene solución.
  \item \medio\ Resolver el sistema $x+y=2$, $x-y=0$, $2x+y=3$ (caso
        $3\times2$).
  \item \medio\ Mostrar que el sistema $x+y=1$, $x-y=3$, $2x+2y=5$ (caso
        $3\times2$) es incompatible.
\end{mwbloque}

\mwbloquehead{Bloque E — Multiplicación de matrices}{46}{34}
\begin{mwbloque}
  \item \facil\ Calcular
        $\mwmat{1&2\\3&4}\mwmat{5&6\\7&8}$.
  \item \facil\ Calcular
        $\mwmat{2&-1\\3&0}\mwmat{1&4\\-2&5}$.
  \item \medio\ Con $A=\mwmat{1&2\\3&4}$ y $B=\mwmat{0&1\\1&0}$, calcular
        $AB$ y $BA$. ¿Son iguales?
  \item \medio\ Calcular
        $\mwmat{1&2&3\\4&5&6}\mwmat{1&0\\0&1\\1&1}$.
  \item \medio\ Calcular
        $\mwmat{1&2\\3&4\\5&6}\mwmat{1&0&1\\0&1&1}$.
  \item \medio\ Indicar cuáles productos están definidos y, en ese caso, su
        orden: (a) $(2\times3)(3\times4)$; (b) $(3\times2)(3\times2)$;
        (c) $(1\times4)(4\times1)$; (d) $(2\times2)(2\times3)$.
  \item \facil\ Verificar que $AI=A$ e $IA=A$ para
        $A=\mwmat{3&-1\\2&5}$.
  \item \medio\ Comprobar que $AB=O$ con
        $A=\mwmat{1&1\\2&2}$ y $B=\mwmat{1&-1\\-1&1}$.
  \item \medio\ Si $A=\mwmat{1&1\\0&1}$, calcular $A^2$, $A^3$ y conjeturar
        $A^n$.
  \item \medio\ Si $A=\mwmat{2&-1\\4&-2}$, calcular $A^2$. ¿Se cumple que
        $A=O$?
  \item \dificil\ Con $A=\mwmat{1&2\\0&1}$ y $B=\mwmat{1&0\\3&1}$, calcular
        $(A+B)^2$ y $A^2+2AB+B^2$. ¿Son iguales? Explicar.
  \item \dificil\ Hallar $a$ y $b$ para que
        $\mwmat{a&b\\1&2}\mwmat{1&1\\1&-1}=\mwmat{3&1\\3&-1}$.
\end{mwbloque}

\mwbloquehead{Bloque F — Matrices equivalentes}{54}{46}
\begin{mwbloque}
  \item \facil\ Aplicar $F_2\leftarrow F_2-2F_1$ a
        $\mwmat{1&2&5\\2&5&12}$.
  \item \facil\ Aplicar $F_1\leftrightarrow F_3$ a
        $\mwmat{1&2\\3&4\\5&6}$.
  \item \facil\ Aplicar $F_2\leftarrow \tfrac13F_2$ a
        $\mwmat{3&6&9\\6&9&12}$.
  \item \medio\ Reducir a forma escalonada
        $\mwmat{1&3&7\\2&4&10}$.
  \item \medio\ ¿Son equivalentes
        $\mwmat{1&2\\2&4}$ y $\mwmat{1&2\\0&0}$? Justificar con las
        operaciones usadas.
  \item \facil\ Indicar cuál de las siguientes \textbf{no} es una operación
        elemental de fila: (a) $F_1\leftrightarrow F_2$;
        (b) $F_1\leftarrow 3F_1$; (c) $F_1\leftarrow 0\cdot F_1$;
        (d) $F_1\leftarrow F_1+2F_2$.
  \item \medio\ Detectar el error: un estudiante aplica
        $F_2\leftarrow F_2-F_1$ solo a las dos primeras columnas de una
        matriz aumentada y deja intacta la columna de constantes. ¿Qué
        ocurre con el sistema?
  \item \medio\ Reducir $\mwmat{0&2&4\\3&6&9}$ a forma escalonada.
\end{mwbloque}

\mwbloquehead{Bloque G — Gauss y Gauss-Jordan}{64}{54}
\begin{mwbloque}
  \item \medio\ Resolver por Gauss:
        $x+y+z=6$, $2x-y+z=6$, $3x+2y-z=5$.
  \item \medio\ Resolver por Gauss:
        $2x+y+z=7$, $x+3y-z=4$, $x+y+2z=9$.
  \item \medio\ Resolver por Gauss:
        $x-y+2z=9$, $2x+y+z=6$, $x+2y+z=3$.
  \item \medio\ Resolver por Gauss-Jordan:
        $x+2y=5$, $2x-y+z=3$, $x+y-z=0$.
  \item \medio\ Resolver por Gauss:
        $x+y+z=3$, $2x+2y+2z=7$, $x-y+z=1$.
  \item \medio\ Resolver por Gauss:
        $x+y+z=6$, $2x-y+z=3$, $3x+2z=9$.
  \item \medio\ Resolver por Gauss:
        $x+y+z=6$, $x+2y+3z=10$, $2x+3y+4z=16$.
  \item \medio\ Resolver por Gauss:
        $2y+z=5$, $x-y+z=0$, $2x+y-z=3$ (atención al pivote nulo).
  \item \medio\ Resolver por Gauss-Jordan:
        $x+y-z=1$, $2x-y+z=2$, $x+2y+2z=5$.
  \item \medio\ Resolver por Gauss-Jordan:
        $3x+y+z=8$, $x+2y-z=3$, $x+y+2z=5$.
\end{mwbloque}

\mwbloquehead{Bloque H — Determinantes $2\times2$ y $3\times3$}{72}{64}
\begin{mwbloque}
  \item \facil\ Calcular $\mwdet{3&4\\2&5}$.
  \item \facil\ Calcular $\mwdet{-2&5\\3&-4}$.
  \item \medio\ Calcular $\mwdet{\tfrac12&2\\[2pt]3&\tfrac13}$.
  \item \medio\ Hallar $x$ si $\mwdet{x&2\\3&x}=10$.
  \item \medio\ Calcular por Sarrus $\mwdet{1&2&3\\4&5&6\\7&8&9}$.
  \item \medio\ Calcular por Sarrus $\mwdet{2&-1&3\\0&4&1\\5&2&-2}$, usando
        la disposición horizontal.
  \item \medio\ Calcular $\mwdet{3&0&2\\1&4&-1\\2&5&3}$ y comprobar el
        resultado con la disposición vertical de Sarrus.
  \item \facil\ Calcular $\mwdet{2&0&0\\5&3&0\\1&4&7}$ sin desarrollar
        cofactores.
\end{mwbloque}

\mwbloquehead{Bloque I — Determinantes $4\times4$ por cofactores}{76}{72}
\begin{mwbloque}
  \item \medio\ Calcular $\mwdet{2&3&1&4\\0&-1&2&5\\0&0&3&6\\0&0&0&4}$.
  \item \dificil\ Calcular
        $\mwdet{1&0&0&0\\2&3&4&5\\1&2&3&4\\0&1&2&3}$ desarrollando por la
        primera fila.
  \item \dificil\ Calcular
        $\mwdet{2&1&0&0\\1&2&1&0\\0&1&2&1\\0&0&1&2}$.
  \item \medio\ Calcular
        $\mwdet{0&1&0&0\\1&0&0&0\\0&0&0&1\\0&0&1&0}$ e interpretar el signo.
\end{mwbloque}

\mwbloquehead{Bloque J — Regla de Cramer}{80}{76}
\begin{mwbloque}
  \item \medio\ Resolver por Cramer:
        $2x+y-z=3$, $x-2y+3z=-4$, $3x+y+2z=9$.
  \item \medio\ Resolver por Cramer:
        $2x-y+z=3$, $x+2y+3z=14$, $3x+y-2z=-1$.
  \item \medio\ Intentar resolver por Cramer:
        $x+2y+3z=4$, $2x+4y+6z=9$, $x+y+z=3$. Explicar qué ocurre.
  \item \medio\ Resolver por Cramer:
        $3x+y-z=4$, $x+2y+z=6$, $2x-y+3z=3$.
\end{mwbloque}

\mwbloquehead{Bloque K — Matriz inversa}{88}{80}
\begin{mwbloque}
  \item \facil\ Verificar que $B=\mwmat{1&-1\\-2&3}$ es la inversa de
        $A=\mwmat{3&1\\2&1}$, calculando $AB$ y $BA$.
  \item \facil\ Hallar la inversa de $\mwmat{4&3\\3&2}$ reduciendo
        $[A\mid I]$.
  \item \medio\ Hallar la inversa de $\mwmat{2&5\\1&3}$ con el atajo de
        $2\times2$ y comprobar el resultado.
  \item \medio\ Intentar hallar la inversa de $\mwmat{1&2\\2&4}$ y explicar
        por qué no existe.
  \item \dificil\ Hallar la inversa de
        $\mwmat{1&2&3\\0&1&4\\5&6&0}$ por Gauss-Jordan.
  \item \dificil\ Hallar la inversa de
        $\mwmat{2&1&1\\1&3&2\\1&0&0}$ por Gauss-Jordan.
  \item \medio\ Resolver con la inversa el sistema $3x+y=5$, $2x+y=3$.
  \item \medio\ Determinar para qué valores de $k$ la matriz
        $\mwmat{k&2\\3&4}$ \textbf{no} tiene inversa.
\end{mwbloque}

% ============================================================
\ifmwclaves
\section*{Clave de respuestas comentada}
% ============================================================

\begin{cuadroinfo}
  \small
  Las respuestas están en el mismo orden del banco. En los ejercicios de
  procedimiento se indica la operación clave; si tu resultado no coincide,
  revisa ese paso antes de repetir todo el cálculo.
\end{cuadroinfo}

\subsection*{Bloque A — Definición, filas, columnas y orden (1--8)}
\begin{enumerate}[label=\textbf{\arabic*.}, start=1, leftmargin=*]
  \item Orden $2\times3$. Filas: $(2,-1,3)$ y $(0,5,-4)$. Columnas:
        $(2,0)$, $(-1,5)$, $(3,-4)$. $a_{12}=-1$, $a_{23}=-4$,
        $a_{21}=0$.
  \item Orden $3\times2$, $6$ elementos. $b_{31}=4$, $b_{22}=3$. Primera
        columna: $(1,-2,4)$.
  \item $C=\mwmat{3&5&7\\4&6&8}$.
  \item $D=\mwmat{0&-1\\1&0\\2&1}$.
  \item $E=\mwmat{3&4\\5&6}$.
  \item $x=4$, $y=5$.
  \item Solo el caso (b): dos matrices $2\times2$ son iguales si coinciden
        todas sus entradas. En (a) y (c) los órdenes son distintos.
  \item (a) $20$; \quad (b) $7$; \quad (c) $n^2$; \quad (d) $12$.
\end{enumerate}

\subsection*{Bloque B — Multiplicación por un escalar (9--16)}
\begin{enumerate}[label=\textbf{\arabic*.}, start=9, leftmargin=*]
  \item $3A=\mwmat{6&-9\\3&12}$.
  \item $-2B=\mwmat{2&-4\\0&-10\\-6&4}$.
  \item $6C=\mwmat{3&-6\\18&4}$.
  \item $2A=\mwmat{4&2\\-2&6}$ y $3B=\mwmat{0&-6\\12&3}$, luego
        $2A-3B=\mwmat{4&8\\-14&3}$.
  \item $2X=\mwmat{4&4\\4&4}$, de donde
        $X=\mwmat{2&2\\2&2}$.
  \item $-\tfrac12A=\mwmat{-\tfrac12&1\\[2pt]-\tfrac32&0}$.
  \item $A+B=\mwmat{0&3\\2&1}$, así que
        $3(A+B)=\mwmat{0&9\\6&3}$. Por otro lado,
        $3A+3B=\mwmat{3&0\\6&-3}+\mwmat{-3&9\\0&6}=\mwmat{0&9\\6&3}$.
        Ambas expresiones coinciden.
  \item $k=-3$, porque $-3(2)=-6$ y $-3(-1)=3$.
\end{enumerate}

\subsection*{Bloque C — Suma y resta de matrices (17--26)}
\begin{enumerate}[label=\textbf{\arabic*.}, start=17, leftmargin=*]
  \item $A+B=\mwmat{6&1\\3&6}$ y $A-B=\mwmat{-4&3\\3&2}$.
  \item $A+B=\mwmat{1&2&-1\\4&1&7}$ y
        $A-B=\mwmat{3&-2&-1\\2&7&3}$.
  \item $A+B=\mwmat{3&0\\1&-3\\5&6}$.
  \item No es posible: la suma exige el mismo orden y $2\times3\neq3\times2$.
        El hecho de que ambas tengan $6$ elementos no alcanza.
  \item $X=\mwmat{4&-1\\5&2}+\mwmat{2&1\\0&3}
        =\mwmat{6&0\\5&5}$.
  \item $a+1=4\Rightarrow a=3$; $b+4=7\Rightarrow b=3$. Las otras dos
        entradas también coinciden: $2+(-2)=0$ y $3+(-3)=0$.
  \item $2A=\mwmat{2&0\\4&-2}$, $2A-B=\mwmat{-1&-1\\5&-4}$ y
        $2A-B+C=\mwmat{-1&3\\6&-4}$.
  \item $A+O=\mwmat{3&-1\\2&0}=A$: la matriz nula no altera la suma.
  \item $A+B=\mwmat{6&1\\3&6}$, $A-B=\mwmat{-4&3\\3&2}$;
        $(A+B)-(A-B)=\mwmat{10&-2\\0&4}=2B$.
  \item $x-1=0\Rightarrow x=1$; $y-3=0\Rightarrow y=3$. Las demás
        entradas ya coinciden.
\end{enumerate}

\subsection*{Bloque D — Sistemas y matriz aumentada (27--34)}
\begin{enumerate}[label=\textbf{\arabic*.}, start=27, leftmargin=*]
  \item $\mwaug{ccc|c}{1&2&-1&3\\2&-1&1&0\\3&1&2&7}$.
  \item Sistema: $x+y+z=6$, $2x-y+z=3$, $x+2y-z=2$. Con $(1,2,3)$:
        $1+2+3=6$ \checkmark, $2-2+3=3$ \checkmark\ y $1+4-3=2$ \checkmark.
        Sí es solución (y es la única).
  \item
        $\mwmat{2&-1&3\\1&4&-1\\3&1&1}\mwmat{x\\y\\z}=\mwmat{5\\2\\8}$.
  \item $A=\mwmat{1&2&3\\2&-1&1\\1&1&-1}$,
        $\mathbf{x}=\mwmat{x\\y\\z}$, $\mathbf{b}=\mwmat{6\\1\\2}$ y
        $\mwaug{ccc|c}{1&2&3&6\\2&-1&1&1\\1&1&-1&2}$.
  \item Aumentada: $\mwaug{ccc|c}{1&1&1&3\\2&-1&1&0}$. Con $z$ libre,
        $x=1-\tfrac23z$ e $y=2-\tfrac13z$. Ejemplos: $z=0$ da
        $(1,2,0)$; $z=3$ da $(-1,1,3)$.
        \emph{Verificación:} $-1+1+3=3$ y $-2-1+3=0$. \checkmark
  \item Las dos ecuaciones tienen los mismos coeficientes pero constantes
        distintas: si $x+y+z$ valiera a la vez $1$ y $2$, se tendría
        $1=2$, lo cual es imposible.
  \item De $x+y=2$ y $x-y=0$ se obtiene $x=1$, $y=1$. La tercera ecuación
        se cumple: $2(1)+1=3$. Solución única: $(1,1)$.
  \item De las dos primeras, $x=2$, $y=-1$. En la tercera:
        $2(2)+2(-1)=2\neq5$. El sistema es incompatible.
\end{enumerate}

\subsection*{Bloque E — Multiplicación de matrices (35--46)}
\begin{enumerate}[label=\textbf{\arabic*.}, start=35, leftmargin=*]
  \item $\mwmat{19&22\\43&50}$.
  \item $\mwmat{4&3\\3&12}$.
  \item $AB=\mwmat{2&1\\4&3}$ y $BA=\mwmat{3&4\\1&2}$. No son iguales:
        la multiplicación no es conmutativa.
  \item $\mwmat{4&5\\10&11}$.
  \item $\mwmat{1&2&3\\3&4&7\\5&6&11}$.
  \item (a) definido, orden $2\times4$; (b) no definido
        ($2\neq3$); (c) definido, orden $1\times1$; (d) definido, orden
        $2\times3$.
  \item $AI=IA=\mwmat{3&-1\\2&5}$: la identidad $2\times2$ deja la matriz
        igual.
  \item $AB=\mwmat{1-1&-1+1\\2-2&-2+2}=\mwmat{0&0\\0&0}$.
  \item $A^2=\mwmat{1&2\\0&1}$, $A^3=\mwmat{1&3\\0&1}$ y, en general,
        $A^n=\mwmat{1&n\\0&1}$.
  \item $A^2=\mwmat{4-4&-2+2\\8-8&-4+4}=\mwmat{0&0\\0&0}$, pero
        $A\neq O$. Un producto de matrices puede anularse sin que ningún
        factor sea nulo.
  \item $A+B=\mwmat{2&2\\3&2}$, entonces
        $(A+B)^2=\mwmat{10&8\\12&10}$. Por otro lado,
        $A^2=\mwmat{1&4\\0&1}$, $B^2=\mwmat{1&0\\6&1}$ y
        $AB=\mwmat{7&2\\3&1}$, así que
        $A^2+2AB+B^2=\mwmat{16&8\\12&4}$. No son iguales porque
        $AB\neq BA$; la fórmula del binomio solo vale si las matrices
        conmutan.
  \item Al multiplicar, la primera fila da $(a+b,\;a-b)$ y la segunda,
        $(3,-1)$. Entonces $a+b=3$ y $a-b=1$, de donde $a=2$, $b=1$.
\end{enumerate}

\subsection*{Bloque F — Matrices equivalentes (47--54)}
\begin{enumerate}[label=\textbf{\arabic*.}, start=47, leftmargin=*]
  \item $\mwmat{1&2&5\\0&1&2}$.
  \item $\mwmat{5&6\\3&4\\1&2}$.
  \item $\mwmat{3&6&9\\2&3&4}$.
  \item $F_2\leftarrow F_2-2F_1$ da $\mwmat{1&3&7\\0&-2&-4}$; luego
        $F_2\leftarrow-\tfrac12F_2$ da $\mwmat{1&3&7\\0&1&2}$.
  \item Sí: aplicando $F_2\leftarrow F_2-2F_1$ a la primera se obtiene
        exactamente la segunda.
  \item La opción (c): multiplicar una fila por $0$ destruye información y
        no es una operación elemental válida.
  \item La operación debe aplicarse a \textbf{toda} la fila. Si la columna
        de constantes queda intacta, la nueva matriz aumentada no representa
        un sistema equivalente y la solución puede cambiar.
  \item $F_1\leftrightarrow F_2$ da $\mwmat{3&6&9\\0&2&4}$;
        $F_1\leftarrow\tfrac13F_1$ da $\mwmat{1&2&3\\0&2&4}$;
        $F_2\leftarrow\tfrac12F_2$ da $\mwmat{1&2&3\\0&1&2}$.
\end{enumerate}

\subsection*{Bloque G — Gauss y Gauss-Jordan (55--64)}
\begin{enumerate}[label=\textbf{\arabic*.}, start=55, leftmargin=*]
  \item Solución $(2,1,3)$. Comprobación: $2+1+3=6$, $4-1+3=6$ y
        $6+2-3=5$.
  \item Solución $(1,2,3)$. Comprobación: $2+2+3=7$, $1+6-3=4$ y
        $1+2+6=9$.
  \item Solución $(2,-1,3)$. Comprobación: $2+1+6=9$, $4-1+3=6$ y
        $2-2+3=3$.
  \item Solución $(1,2,3)$. Gauss-Jordan deja
        $\mwaug{ccc|c}{1&0&0&1\\0&1&0&2\\0&0&1&3}$.
  \item La segunda ecuación es el doble de la primera, pero su constante no:
        pide $2x+2y+2z=7$ cuando debería pedir $6$. Al reducir aparece
        $[0,0,0\mid1]$: \textbf{sin solución}.
  \item La tercera ecuación es la suma de las dos primeras
        ($3x+2z=9$). Queda un sistema $2\times3$ con $z=t$ libre:
        \[
          x=3-\tfrac23t,\qquad y=3-\tfrac13t,\qquad z=t.
        \]
        Infinitas soluciones.
  \item La tercera ecuación es también la suma de las dos primeras; queda
        un sistema $2\times3$. Con $z=t$:
        \[
          x=2+t,\qquad y=4-2t,\qquad z=t.
        \]
        Infinitas soluciones.
  \item El pivote de la primera columna es $0$: hay que intercambiar
        $F_1\leftrightarrow F_2$. Solución $(1,2,1)$.
  \item Solución $(1,1,1)$. Comprobación: $1+1-1=1$, $2-1+1=2$ y
        $1+2+2=5$.
  \item Solución $(2,1,1)$. Comprobación: $6+1+1=8$, $2+2-1=3$ y
        $2+1+2=5$.
\end{enumerate}

\subsection*{Bloque H — Determinantes $2\times2$ y $3\times3$ (65--72)}
\begin{enumerate}[label=\textbf{\arabic*.}, start=65, leftmargin=*]
  \item $15-8=7$.
  \item $8-15=-7$.
  \item $\tfrac16-6=-\tfrac{35}{6}$.
  \item $x^2-6=10\Rightarrow x^2=16\Rightarrow x=\pm4$.
  \item Por Sarrus: $225-225=0$. La tercera fila es combinación de las
        anteriores.
  \item $(-16-5+0)-(60+4+0)=-21-64=-85$.
  \item $3(12+5)-0+2(5-8)=51-6=45$.
  \item Es triangular inferior: el determinante es el producto de la
        diagonal, $2\cdot3\cdot7=42$.
\end{enumerate}

\subsection*{Bloque I — Determinantes $4\times4$ por cofactores (73--76)}
\begin{enumerate}[label=\textbf{\arabic*.}, start=73, leftmargin=*]
  \item Es triangular superior: $2\cdot(-1)\cdot3\cdot4=-24$.
  \item Desarrollando por la primera fila:
        $\det=1\cdot\mwdet{3&4&5\\2&3&4\\1&2&3}$.
        Ese menor vale $3(9-8)-4(6-4)+5(4-3)=3-8+5=0$.
        Por lo tanto, $\det=0$.
  \item Desarrollo por la primera fila:
        \[
          \det=2\mwdet{2&1&0\\1&2&1\\0&1&2}
               -1\mwdet{1&1&0\\0&2&1\\0&1&2}.
        \]
        El primer menor vale $2(4-1)-1(2-0)=4$; el segundo,
        $1(4-1)-1(0)=3$. Entonces
        $\det=2(4)-3=5$.
  \item Intercambiando $F_1\leftrightarrow F_2$ y $F_3\leftrightarrow F_4$
        se obtiene $I_4$. Cada intercambio cambia el signo, así que
        $\det=(-1)^2=1$.
\end{enumerate}

\subsection*{Bloque J — Regla de Cramer (77--80)}
\begin{enumerate}[label=\textbf{\arabic*.}, start=77, leftmargin=*]
  \item $\det(A)=-14$; $\det(A_x)=0$, $\det(A_y)=-70$ y
        $\det(A_z)=-28$. Entonces $x=0$, $y=5$, $z=2$.
  \item $\det(A)=-30$; $\det(A_x)=-30$, $\det(A_y)=-60$ y
        $\det(A_z)=-90$. Solución $(1,2,3)$.
  \item $\det(A)=0$: Cramer no se puede aplicar porque habría que dividir
        entre cero. Además la segunda ecuación es el doble de la primera
        ($2x+4y+6z=8$) pero su constante es $9$, así que el sistema es
        \textbf{incompatible}.
  \item $\det(A)=25$; $\det(A_x)=25$, $\det(A_y)=50$ y $\det(A_z)=25$.
        Solución $(1,2,1)$.
\end{enumerate}

\subsection*{Bloque K — Matriz inversa (81--88)}
\begin{enumerate}[label=\textbf{\arabic*.}, start=81, leftmargin=*]
  \item $AB=\mwmat{3(1)+1(-2) & 3(-1)+1(3)\\-2+2 & -2+3}
        =\mwmat{1&0\\0&1}$ y $BA=I_2$ también. Luego $B=A^{-1}$.
  \item Por Gauss-Jordan se llega a
        $\mwaug{cc|cc}{1&0&-2&3\\0&1&3&-4}$, es decir
        $A^{-1}=\mwmat{-2&3\\3&-4}$ ($\det A=-1$).
  \item $\det(A)=6-5=1$, así que
        $A^{-1}=\dfrac{1}{1}\mwmat{3&-5\\-1&2}=\mwmat{3&-5\\-1&2}$.
        Comprobación: $2(3)+5(-1)=1$ y $2(-5)+5(2)=0$. \checkmark
  \item Al reducir aparece la fila $[0,0\mid-2,1]$, con la mitad izquierda
        nula: no se puede llegar a $I_2$. Además
        $\det(A)=4-4=0$. La matriz es \textbf{singular}.
  \item $A^{-1}=\mwmat{-24&18&5\\20&-15&-4\\-5&4&1}$, según el ejemplo 3 de
        la sección 9. Verificación: primera fila de $A$ por primera columna de
        $A^{-1}$: $-24+40-15=1$.
  \item $\det(A)=-1$ y
        $A^{-1}=\mwmat{0&0&1\\-2&1&3\\3&-1&-5}$. Comprobación:
        primera fila de $A$ por primera columna de $A^{-1}$:
        $2(0)+1(-2)+1(3)=1$. \checkmark
  \item $A^{-1}=\mwmat{1&-1\\-2&3}$, luego
        $\mathbf{x}=A^{-1}\mathbf{b}
        =\mwmat{1&-1\\-2&3}\mwmat{5\\3}=\mwmat{2\\-1}$.
        Solución $(2,-1)$.
  \item $\det(A)=4k-6$. No hay inversa cuando $4k-6=0$, es decir
        $k=\tfrac32$. Para cualquier otro valor, la inversa existe.
\end{enumerate}

\fi
\end{document}
